\chapter{Results}\label{ch:results}

The results follow the order of the study: D3IL-avoiding, the benchmark of \ac{DPCC}, first, then
the two environments built for this thesis, D3IL-aligning and the UAV scenes. \autoref{sec:res:summary} brings the three together and
answers, per environment, how the generative models and the projection methods compare. Each
environment is reported in those two parts and closes with the supported model--projection combination.
Each table states its backbone and its evaluation protocol.

\section{Obstacle Avoidance}\label{sec:res:state}

Every comparison in this section is made at the protocol of \ac{DPCC} --- five training seeds and two
episodes per seed and geometry, ten episodes per geometry (\autoref{sec:setup:protocol:avoiding}) ---
except where a table states otherwise: \autoref{tab:avoiding-projectors} runs on four training seeds at
$\nfe=3$ and on one at $\nfe=10$, one of its cells on twenty episodes.
Results are averaged per geometry and then across the three geometries. At every control step the planner generates
$B=4$ candidate plans and executes one of them, and every model is shown under each of the three rules
that make that choice (\autoref{sec:method:selection}): random selection ($r$), cumulative projection
cost ($c$) and temporal consistency ($t$). The models are then compared at their best rule.

The generative-model comparison first reads the raw outputs with projection switched off. All results
with a projector in the loop are then reported together under projection methods, including the
comparison of the four models through the fixed per-step projector of \ac{DPCC} and the comparison
between per-step and endpoint projection. The section closes with the supported model--projection
combinations.

\subsection{Generative Models}\label{sec:res:avoiding:models}

The four models of \autoref{sec:setup:baselines} are compared, on the shared temporal U-Net: FM,
MeanFM, CI-MeanFM and the diffusion model of \ac{DPCC}. One property of the step budget decides how the
results below are read. $\nfe$ counts network evaluations in both families but is not the same quantity
in them (\autoref{tab:eval}): for the flow-based models it is the number of ODE solver steps, chosen at
inference, and for the baseline the number of denoising steps, fixed when its noise schedule is
discretised during training --- so every diffusion budget is a separately trained checkpoint and every
flow-based budget is the same checkpoint read again. This subsection reads the plans executed without
projection; projected results begin in \autoref{sec:res:avoiding:projection}.

The demonstrations teach the route to the goal. With projection switched off, each of the three
flow-based models reaches it in all thirty episodes at both reported budgets; the diffusion model
reaches it in 28 of thirty at one denoising step and 29 of thirty at twenty
(\autoref{tab:avoiding-raw-models}). Reaching the goal is therefore distinct from satisfying the
test-time constraints: no unprojected row exceeds $0.133$ success with constraint satisfaction.

\subsubsection{Raw Model Outputs}\label{sec:res:avoiding:raw}

With the projector switched off, \autoref{tab:avoiding-raw-models} measures goal arrival, constraint satisfaction and generation time on the same cells as \autoref{tab:avoiding-dpcc-protocol}.

\begin{table}[htpb]
  \caption[The four models without projection]{D3IL-avoiding with the projection switched off: the
  cells of \autoref{tab:avoiding-dpcc-protocol} at the same protocol, five training seeds $\times$ two
  episodes per geometry, averaged per geometry and then over the three; tightened constraints, four
  candidate plans, temporal U-Net of $4.0$\,M parameters. Success, S\&C and violating steps per episode as
  in \autoref{sec:setup:metrics:avoiding}; each cell mean $\pm$ standard deviation over the seeds. No rule
  column: the evaluation records one unprojected arm per cell. The budgets $\nfe=5$ and $10$ were
  evaluated as well and are not printed. \baselinemark\ the baseline of record at its own budget.}
  \label{tab:avoiding-raw-models}
  \centering
  \footnotesize
  \begin{tabularx}{\linewidth}{@{}L r rrrrr@{}}
    \toprule
    Model & $\nfe$ & Success & S\&C & \shortstack{Violating\\steps} & Steps & ms/step \\
    \midrule
    MeanFM    & 1  & $1.000 \pm 0.000$ & $0.033 \pm 0.075$ & $15.5 \pm 2.1$ & $63.4 \pm 5.6$ & $10.3 \pm 0.4$ \\
    MeanFM    & 2  & $1.000 \pm 0.000$ & $0.033 \pm 0.075$ & $15.5 \pm 0.4$ & $62.7 \pm 1.8$ & $19.3 \pm 0.2$ \\
    \midrule
    CI-MeanFM, $\alpha_{\mathrm{end}}=0.2$ & 1 & $1.000 \pm 0.000$ & $0.033 \pm 0.075$ & $15.6 \pm 0.7$ & $60.5 \pm 1.5$ & $9.8 \pm 0.2$ \\
    CI-MeanFM, $\alpha_{\mathrm{end}}=0.2$ & 2 & $1.000 \pm 0.000$ & $0.033 \pm 0.075$ & $15.7 \pm 0.7$ & $62.9 \pm 1.2$ & $18.7 \pm 0.3$ \\
    \midrule
    FM        & 1  & $1.000 \pm 0.000$ & $0.100 \pm 0.091$ & $16.5 \pm 3.4$ & $68.6 \pm 5.9$ & $9.0 \pm 0.0$ \\
    FM        & 2  & $1.000 \pm 0.000$ & $0.100 \pm 0.091$ & $15.6 \pm 2.6$ & $61.7 \pm 2.2$ & $17.8 \pm 0.0$ \\
    FM        & 20 & $1.000 \pm 0.000$ & $0.133 \pm 0.075$ & $16.5 \pm 2.1$ & $70.2 \pm 6.3$ & $183.5 \pm 0.7$ \\
    \midrule
    Diffusion & 1  & $0.933 \pm 0.149$ & $0.000 \pm 0.000$ & $19.0 \pm 2.9$ & $68.3 \pm 16.1$ & $9.4 \pm 0.1$ \\
    Diffusion & 2  & $1.000 \pm 0.000$ & $0.000 \pm 0.000$ & $20.4 \pm 2.4$ & $61.7 \pm 8.3$ & $18.3 \pm 0.1$ \\
    \mainconfig{\baselinemark\ Diffusion} & \mainconfig{20} & \mainconfig{$0.967 \pm 0.075$} & \mainconfig{$0.100 \pm 0.091$} & \mainconfig{$17.9 \pm 2.5$} & \mainconfig{$67.8 \pm 16.8$} & \mainconfig{$179.3 \pm 1.2$}  \\
    \bottomrule
  \end{tabularx}
\end{table}

Read down the success column first. Every flow-based model reaches the goal in every one of its thirty
episodes at every budget in the table, and the baseline in $0.933$ of them at one denoising step and
$0.967$ at twenty. Producing a plan that arrives is not what separates these models on this benchmark.

Read the next column and none of them arrives legally. The highest success with constraint satisfaction
anywhere in the table is $0.133$, instantaneous-velocity matching at twenty evaluations; the two
average-velocity models read $0.033$ at the budgets they are operated at, instantaneous-velocity
matching $0.100$, the baseline $0.100$ at its own budget and $0.000$ at one denoising step. Between
$15.5$ and $20.4$ control steps of an average episode violate a constraint, and raising the budget does
not work that number down: the intermediate budgets of five and ten, which the table does not print,
hold it between $14.2$ and $19.1$, and the best and the worst of those counts both come at a budget of ten.
The demonstrations all four models learn from cross the
same constraint sets (\autoref{tab:avoiding-geometries}), so a plan that respects them is not something
any of them could have imitated. On this benchmark constraint satisfaction is the projector's work, for
the flow-based objectives exactly as for the diffusion baseline.

The last column prices generation alone. One network evaluation of a flow-based model costs $9.0$ to
$10.3$\,ms per control step and twenty denoising steps of the baseline $179.3$\,ms. Against the projected
table the remainder is the projector: about $8$ of the $18.3$\,ms analytic average-velocity matching
takes at one network evaluation, and about $374$ of the baseline's $553.4$\,ms at its own twenty
denoising steps. The projector runs on
$\lceil \eta\nfe \rceil$ sampling steps (\autoref{sec:setup:protocol:guiding}), which is one at the
flow-based operating budget and eleven at the baseline's, whose gate includes the threshold step, so the
budget is paid twice --- once in the network and again in the solver.

\FloatBarrier

Those three columns score what the models produce. \autoref{fig:raw-plans} shows it. The raw plan is to a
planning model what a generated image is to an image model: the fidelity of the sample itself, before any
use is made of it \parencite{ho2020denoising,lipman2023flow,geng2025meanflow}, and the columns above are
that fidelity scored on the task. The one thing to be clear about before reading the figure is that it is
not the executed path --- the executed path is not drawn in it
at all. At every control step the policy returns a plan tensor of shape $(N,8,6)$
(\autoref{tab:train}): $N=4$ candidate plans, each $H=8$ waypoints long, each waypoint carrying a planar
position increment together with the commanded and the measured planar position. What the figure draws,
per candidate, is the \emph{measured} planar position channels of those eight waypoints --- one
eight-point curve --- and the green dot marks the first of the eight, the state the plan is anchored at.
Drawing the fan at every control step would fill the panel with a solid block, so it is drawn at every
fourth; the stored tensor keeps all of them.

Of the eight waypoints exactly one is ever carried out. The selection rule picks one of the four
candidates (\autoref{sec:method:selection}), its first action is applied as a Cartesian increment, and the
other seven waypoints are thrown away and regenerated at the next control step, because the planner
replans every step (\autoref{tab:eval}). So the eight-step plan is never executed as drawn --- and that is
what makes it worth drawing. It is the network's output before the projector repairs it and before the
controller smooths it, the quantity that for an image generative model would be the fidelity of a sample.
Every other figure and table in this section reads the models through something downstream; this one reads
them directly.

\begin{figure}[!htbp]
  \centering
  \setlength{\tabcolsep}{2pt}
  \begin{tabular}{@{}c cccc@{}}
    & {\footnotesize\bfseries MeanFM $\downarrow$}
    & {\footnotesize\bfseries CI-MeanFM $\downarrow$}
    & {\footnotesize FM}
    & {\footnotesize Diffusion} \\[0.2em]
    {\footnotesize $\nfe=1$} &
    \begin{minipage}[c]{0.21\linewidth}\centering
    \includegraphics[width=\linewidth]{fig_raw_plans_meanflow_K1}
    \end{minipage} &
    \begin{minipage}[c]{0.21\linewidth}\centering
    \includegraphics[width=\linewidth]{fig_raw_plans_af_K1}
    \end{minipage} &
    \begin{minipage}[c]{0.21\linewidth}\centering
    \includegraphics[width=\linewidth]{fig_raw_plans_fm_K1}
    \end{minipage} &
    \begin{minipage}[c]{0.21\linewidth}\centering
    \includegraphics[width=\linewidth]{fig_raw_plans_diffusion_K1}
    \end{minipage} \\
    {\footnotesize $\nfe=2$} &
    \begin{minipage}[c]{0.21\linewidth}\centering
    \includegraphics[width=\linewidth]{fig_raw_plans_meanflow_K2}
    \end{minipage} &
    \begin{minipage}[c]{0.21\linewidth}\centering
    \includegraphics[width=\linewidth]{fig_raw_plans_af_K2}
    \end{minipage} &
    \begin{minipage}[c]{0.21\linewidth}\centering
    \includegraphics[width=\linewidth]{fig_raw_plans_fm_K2}
    \end{minipage} &
    \begin{minipage}[c]{0.21\linewidth}\centering
    \includegraphics[width=\linewidth]{fig_raw_plans_diffusion_K2}
    \end{minipage} \\
    {\footnotesize $\nfe=20$} &
    \begin{minipage}[c]{0.21\linewidth}\centering{\footnotesize ---}\end{minipage} &
    \begin{minipage}[c]{0.21\linewidth}\centering{\footnotesize ---}\end{minipage} &
    \begin{minipage}[c]{0.21\linewidth}\centering{\footnotesize ---}\end{minipage} &
    \begin{minipage}[c]{0.21\linewidth}\centering
    \includegraphics[width=\linewidth]{fig_raw_plans_diffusion_K20}
    \end{minipage} \\
  \end{tabular}
  \caption[Plans without projection, four models across the step budget]{Plans without projection, one episode
  per panel: the plans the model emitted at every fourth control step, all four candidates, each drawn as
  the measured planar position of its eight waypoints, green at the waypoint the plan is anchored at; the
  executed path is not drawn. Columns: MeanFM, CI-MeanFM, FM and the diffusion model of \ac{DPCC}; rows:
  the step budget; \emph{both-hard}, one training seed, temporal U-Net of $4.0$\,M parameters. The arrows
  mark the two average-velocity columns; the last row is the baseline at its trained budget, where the
  flow-based models are not run (dashes). \autoref{tab:avoiding-raw-models} aggregates the outcomes.}
  \label{fig:raw-plans}
\end{figure}

Read that way, the models separate (\autoref{fig:raw-plans}). At one network evaluation the successive
plans of analytic average-velocity matching form a narrow ribbon: consecutive fans overlap, and each of
the eight-point curves is a bounded arc towards the goal. The diffusion baseline at the same budget emits
plans that jump across the workspace from one control step to the next.

That difference does not appear in the motion, and it is worth saying why before the panels are used for
anything. Two things sit between a plan and the arm (\autoref{sec:setup:platforms}): only the
\emph{first} action of the plan is executed, and that action is a Cartesian increment reaching the joints
through inverse kinematics and a joint-space PD controller. A path assembled one first-action at a time,
through an integrator, is smooth whether or not the plans behind it were. The panels are therefore the
only place the incoherence is visible at all. The analytic average-velocity matching panel at two
evaluations shows the ribbon surviving a second evaluation.

The two remaining objectives, in the middle two columns, fall between those extremes and in the order
the projected results below give. Instantaneous-velocity matching throws a wide, jagged fan at both budgets, and much of its length
lies inside the forbidden halfspace on the left of the panel rather than in the free corridor.
Consistency-interpolated average-velocity matching is noisy at one evaluation and markedly tighter at two,
where its plans thread the obstacle field in a band comparable to the analytic objective's. Neither is as
tight as analytic average-velocity matching at the same budget, and all three are far from the near-scribble
of the baseline at one evaluation.

The baseline's small-budget panels are two models, one trained at each budget. At $\nfe=2$ its plans
form a single coherent bundle from start to goal, and that bundle runs along the wrong side of the
constraint. At its own twenty steps, in the last row, the baseline draws a single tight bundle from the
start state to the goal line, the upper half of which lies inside the left halfspace.

In the cell these panels are drawn from --- one training seed, \emph{both-hard}, two episodes, no
projection --- the violating control steps per episode are $6.0$ for analytic average-velocity matching at one evaluation and $12.0$ at two, and for the
diffusion baseline $28.0$ at one, $21.5$ at two and $18.5$ at twenty.

Each panel is one episode. The panels share the geometry and the training seed and come from separate
evaluations. A flow-based column is one checkpoint read at both budgets; the three diffusion panels are
three separately trained models, one per budget, since a diffusion model's step count is fixed at
training. The last row is the baseline at twenty denoising steps; the flow-based models are operated at
one and two evaluations and have no panel there.

\subsection{Projection Methods}\label{sec:res:avoiding:projection}

The remaining results all include projection. The first comparison holds the per-step projection of
\ac{DPCC} fixed and varies the generative model. The later comparison holds the generative model fixed
and varies the projection method.

\subsubsection{At the Protocol of DPCC}\label{sec:res:avoiding:dpccprotocol}\label{sec:res:fewstep}

This pass puts the projector back and holds it fixed at the per-step projection of \ac{DPCC}, so that
the generative model is the only thing that varies; the two projection methods are compared against
each other in \autoref{sec:res:avoiding:projectors}. Every number from here to the end of the
subsection is an executed-episode result of that pipeline.

Two marks recur in this table and in the ones after it. The \textbf{operating point}, \selectedmark, is
the configuration this thesis carries forward on the benchmark --- the cheapest one that holds its
success, not the highest number in a column. The \textbf{baseline of record}, \baselinemark, is the
diffusion model of \ac{DPCC} at its own budget of twenty under its own best selection rule, and it is
what every claim of this chapter is measured against. Every row of the table carries all five training
seeds and all three geometries: the budgets are $\nfe\in\{1,2,5,10\}$ for analytic average-velocity
matching, $\{1,2\}$ for the consistency-interpolated model, $\{1,2,20\}$ for instantaneous-velocity
matching and $\{1,2,10,20\}$ for the diffusion baseline, whose step count is fixed at training.

The baseline reproduces its published result (\autoref{tab:avoiding-dpcc-protocol},
\autoref{tab:target}): under cumulative projection cost it reaches a success rate with constraint
satisfaction of $1.000$ in $70.1$ control steps, at $553.4$\,ms per control step. Temporal consistency
reaches the same rate in more steps and random selection falls short of it.

At two denoising steps under cumulative projection cost the baseline also reaches $1.000$, in $61.2$
control steps at $211.4$\,ms, on five seeds and three geometries: the same constraint satisfaction in
fewer control steps at $38\,\%$ of the time. At ten steps under temporal consistency it reaches $1.000$
in $68.7$ control steps at $321.7$\,ms. The baseline of record is
the configuration \ac{DPCC} reports, at twenty steps, and the cost factors below are stated against it;
against the two-step diffusion configuration the $30\times$ below is $12\times$.

That two-step time is six times the one-step time for twice the denoising steps, and where the
difference comes from is set out in \autoref{sec:res:avoiding:caveat:noise}.

Two flow-based models clear that row at a single network evaluation, on all three of its axes at once.
Instantaneous-velocity matching reaches $1.000$ in $67.0$ control steps and consistency-interpolated
average-velocity matching $1.000$ in $59.2$, the fewest in the table, each at about $18$\,ms per control
step --- a thirtieth of the baseline's, on the same temporal U-Net of $4.0$\,M parameters. Neither gives
up success with constraint satisfaction, and both finish in fewer control steps. Here a flow-based model
is ahead of the baseline on every reported axis at once rather than trading one against another.

Analytic average-velocity matching, the model this benchmark is operated with, reaches $0.967$ in $58.6$
control steps at $18.3$\,ms --- against the baseline's best rule, one episode of thirty fewer reaching
the goal within the constraints, $11.5$ fewer control steps, and a thirtieth of the time. What moves that
result is the selection rule: under cumulative projection cost the same model drops to $0.933$ and stalls
at $72.4$ or $97.2$ control steps, the
behaviour described in \autoref{sec:res:ablations}. The three models differ in how exposed they are to
that choice. Instantaneous-velocity matching is the only one that holds $1.000$ under all three rules at
both budgets; consistency-interpolated average-velocity matching holds it under temporal consistency at
both, but falls to $0.900$ under random selection at two evaluations and stalls under cumulative
projection cost in the same way.

\begin{table}[htpb]
  \caption[The generative models at the protocol of DPCC]{D3IL-avoiding at the protocol of \ac{DPCC}:
  five training seeds $\times$ two episodes per geometry, averaged per geometry and then over the three;
  tightened constraints, four candidate plans, temporal U-Net of $4.0$\,M parameters. Rules: $r$ random,
  $c$ cumulative projection cost, $t$ temporal consistency. Columns as in
  \autoref{sec:setup:metrics:avoiding}, each mean $\pm$ standard deviation over the seeds. The budgets
  $\nfe=5$ and $10$ were evaluated as well and are not printed. \selectedmark\ the operating
  point, \baselinemark\ the baseline of record, bold the cells that clear it on all three columns.}
  \label{tab:avoiding-dpcc-protocol}
  \centering
  \footnotesize
  \begin{tabularx}{\linewidth}{@{}L r c rrr@{}}
    \toprule
    Model & $\nfe$ & Rule & S\&C & Steps & ms/step \\
    \midrule
    MeanFM     & 1 & $r$ & $0.900 \pm 0.149$ & $64.3 \pm 2.6$ & $19.2 \pm 1.6$ \\
                     &    & $c$ & $0.933 \pm 0.091$ & $72.4 \pm 8.4$ & $18.3 \pm 0.5$ \\
                     &    & \mainconfig{\selectedmark\ $t$} & \mainconfig{$0.967 \pm 0.075$} & \mainconfig{$58.6 \pm 2.2$} & \mainconfig{$18.3 \pm 0.2$}  \\
    MeanFM     & 2 & $r$ & $0.967 \pm 0.075$ & $65.1 \pm 2.6$ & $27.6 \pm 0.5$ \\
                     &    & $c$ & $0.933 \pm 0.091$ & $97.2 \pm 4.8$ & $27.1 \pm 0.2$ \\
                     &    & $t$ & $0.967 \pm 0.075$ & $59.4 \pm 2.0$ & $27.7 \pm 0.5$ \\
    \midrule
    CI-MeanFM, $\alpha_{\mathrm{end}}=0.2$ & 1 & $r$ & $1.000 \pm 0.000$ & $61.5 \pm 0.8$ & $17.9 \pm 0.2$ \\
                     &    & $c$ & $0.967 \pm 0.075$ & $68.4 \pm 2.3$ & $18.0 \pm 0.6$ \\
                     &    & \textbf{$t$} & \textbf{$1.000 \pm 0.000$} & \textbf{$59.2 \pm 1.2$} & \textbf{$18.1 \pm 0.2$} \\
    CI-MeanFM, $\alpha_{\mathrm{end}}=0.2$ & 2 & $r$ & $0.900 \pm 0.149$ & $63.3 \pm 1.3$ & $27.0 \pm 0.3$ \\
                     &    & $c$ & $0.933 \pm 0.091$ & $89.7 \pm 8.2$ & $26.6 \pm 0.2$ \\
                     &    & $t$ & $1.000 \pm 0.000$ & $60.1 \pm 3.4$ & $27.0 \pm 0.2$ \\
    \midrule
    FM         & 1 & $r$ & $1.000 \pm 0.000$ & $67.9 \pm 1.2$ & $17.4 \pm 0.1$ \\
                     &    & \textbf{$c$} & \textbf{$1.000 \pm 0.000$} & \textbf{$67.0 \pm 0.8$} & \textbf{$17.3 \pm 0.1$} \\
                     &    & $t$ & $1.000 \pm 0.000$ & $68.2 \pm 1.7$ & $17.7 \pm 0.6$ \\
    FM         & 2 & $r$ & $1.000 \pm 0.000$ & $65.7 \pm 0.8$ & $26.2 \pm 0.9$ \\
                     &    & $c$ & $1.000 \pm 0.000$ & $67.0 \pm 3.8$ & $25.5 \pm 0.2$ \\
                     &    & $t$ & $1.000 \pm 0.000$ & $63.8 \pm 1.9$ & $25.9 \pm 0.2$ \\
    FM         & 20 & $r$ & $0.867 \pm 0.075$ & $73.3 \pm 7.9$ & $852.6 \pm 624.9$ \\
                     &    & $c$ & $1.000 \pm 0.000$ & $63.2 \pm 1.5$ & $476.7 \pm 26.0$ \\
                     &    & $t$ & $1.000 \pm 0.000$ & $63.7 \pm 2.2$ & $486.5 \pm 41.4$ \\
    \midrule
    Diffusion  & 1 & $r$ & $0.600 \pm 0.303$ & $67.7 \pm 10.1$ & $33.9 \pm 8.6$ \\
                     &    & $c$ & $0.667 \pm 0.391$ & $72.1 \pm 17.6$ & $33.3 \pm 14.5$ \\
                     &    & $t$ & $0.600 \pm 0.190$ & $73.5 \pm 14.1$ & $33.3 \pm 13.2$ \\
    Diffusion  & 2 & \textbf{$r$} & \textbf{$1.000 \pm 0.000$} & \textbf{$65.1 \pm 1.4$} & \textbf{$191.1 \pm 26.5$} \\
                     &    & \textbf{$c$} & \textbf{$1.000 \pm 0.000$} & \textbf{$61.2 \pm 1.2$} & \textbf{$211.4 \pm 18.8$} \\
                     &    & $t$ & $0.933 \pm 0.091$ & $60.9 \pm 1.5$ & $225.9 \pm 48.5$ \\
    Diffusion  & 20 & $r$ & $0.967 \pm 0.075$ & $74.7 \pm 5.8$ & $573.2 \pm 22.2$ \\
                     &    & \mainconfig{\baselinemark\ $c$} & \mainconfig{$1.000 \pm 0.000$} & \mainconfig{$70.1 \pm 8.2$} & \mainconfig{$553.4 \pm 28.0$}  \\
                     &    & $t$ & $1.000 \pm 0.000$ & $76.1 \pm 15.7$ & $563.0 \pm 25.4$ \\
    \bottomrule
  \end{tabularx}
\end{table}

\subsubsection{Candidate Plans and Selection Rules}\label{sec:res:ablations}

With a single candidate plan the three selection rules return the same plan: success and control steps
are identical in all fifteen paired seed--geometry combinations, for the diffusion baseline, analytic average-velocity matching,
instantaneous-velocity matching and consistency-interpolated average-velocity matching. With four candidates, generation time changes by at most six per
cent, because the four plans are one batched network evaluation, while projection time grows by a factor
of $3.2$ to $4.9$, because the programs are solved one after another on the CPU. Solving them in parallel
would bring analytic average-velocity matching at $\nfe=1$ from $18.1$ to about $11.7$\,ms per control step with all four
candidates kept; this is a prediction from two measured points.

Cumulative projection cost slows analytic average-velocity matching down on this task. In
\autoref{tab:avoiding-dpcc-protocol} it takes $72.4$ control steps at $\nfe=1$ and $97.2$ at $\nfe=2$,
against $58.6$ and $59.4$ under temporal consistency, and reaches $0.933$ success with constraint
satisfaction against $0.967$: the rule occasionally selects a plan that stalls, and a stalled plan costs
steps without buying constraint satisfaction. The same happens to consistency-interpolated
average-velocity matching, which goes from $60.1$ steps under temporal consistency to $89.7$ under
cumulative projection cost at $\nfe=2$, and it does not happen to the diffusion baseline, for which
cumulative projection cost is the best of the three. Temporal consistency is therefore the selection rule
the flow-based models are operated with on this task.

\subsubsection{The Cost Frontier}\label{sec:res:avoiding:pareto}

At equal success, two costs remain: how many control steps an episode takes, and how long one action takes
to compute. They are not exchangeable --- a plan that arrives in fewer steps but costs half a second per
step is not obviously better than the reverse --- so this section does not add them into one score. It
places every configuration on both axes at once and asks which of them no other configuration improves on in both
(\autoref{fig:avoiding-tradeoff}).

A configuration is \emph{non-dominated} when no other configuration in the comparison is at the same time
faster per action and shorter in control steps. Only configurations that keep the task solved take part:
a point enters the comparison when its success with constraint satisfaction is within $0.05$ of the best
value in its panel, so that a fast configuration that falls further short of the best cannot appear as a
favourable trade; at thirty episodes the band admits $29/30$ beside $30/30$, and where that matters the text
says so. Reading the figure, colour identifies the model, marker shape the selection rule, the
number beside a point its step budget, and a ring marks each non-dominated point; down and to the left is
better on both axes, which the boxed arrow in the legend repeats. The step axis is broken at $80$: every
configuration that solves the task lies below that line, and the few above it --- the hollow $\nfe=2$
points and one diffusion cell --- are drawn on a compressed scale so that the solved band is not squeezed
into a third of the panel. The figure and
\autoref{tab:avoiding-dpcc-protocol} draw on the same cells: every budget plotted here is a row of that
table, on five training seeds and three geometries, under both selection rules, and the frontier is taken
over all of them, so the ringed point of a model is that model at its better rule.

\begin{figure}[htpb]
  \centering
  \includegraphics[width=\linewidth]{fig_avoiding_tradeoff}
  \caption[Control steps against planner time per control step]{Control steps against the time to
  compute one action (generation, projection and selection), every step budget of
  \autoref{tab:avoiding-dpcc-protocol} under both selection rules: circles cumulative projection cost,
  squares temporal consistency; colour the model, the number beside a point its step budget; a hollow
  point is more than $0.05$ below its panel's best success with constraint satisfaction and stays out of
  the comparison; a ring a non-dominated point, the boxed arrow the direction in which both axes improve. Three panels are the
  constraint geometries; the fourth averages within a geometry and then across the three. The step axis
  is broken at $80$. Five training seeds $\times$ two episodes per geometry, tightened constraints,
  temporal U-Net of $4.0$\,M parameters.}
  \label{fig:avoiding-tradeoff}
\end{figure}

The horizontal axis is planner time, not the full executed-step time. Full executed-step timing is
available for the separate UAV controller comparison in \autoref{tab:uav-controller}; it was not
recorded for the matched avoiding configurations plotted here.

The frontier is occupied by the flow-based models at their smallest budgets, and the three of them at
the lower left reach the baseline's control-step count at one twentieth to one thirtieth of its time per
action. Consistency-interpolated average-velocity matching at one network evaluation sits at the lower
left, $59.2$ control steps at $18.1$\,ms, with analytic average-velocity matching beside it at $58.6$
and $18.3$; no configuration of any model is both faster and shorter than either.
Instantaneous-velocity matching at one and two evaluations lies a little longer in steps at a similar
time. Every diffusion configuration lies to the right by more than an order of magnitude in time, and
the baseline's own operating point, twenty evaluations with cumulative projection cost, is neither the
fastest nor the shortest point in the figure.
What the frontier shows is therefore not a single best configuration but a region --- one or two network
evaluations of a flow-based model --- inside which the remaining choices trade a control step against a
millisecond.

\subsubsection{Caveat: What the Projector Is Handed}\label{sec:res:avoiding:caveat:noise}

One number in \autoref{tab:avoiding-dpcc-protocol} does not scale the way the others do, and it says
something about the two sampler families that the table itself cannot. The diffusion baseline's own cost
doubles from one denoising step to two ($9.4$ to $18.3$\,ms, \autoref{tab:avoiding-raw-models}), as it
should; its projected cost rises six-fold, from $33$ to $191$--$226$\,ms. The projection is the same
program, on the same four candidate plans, against the same tightened constraints.

What changes is the state the projector is handed. The baseline projects at every denoising step with
$t \le 0.5\,\nfe$, which at two steps is both of them: the final step, as at every budget, and the first
--- a plan one denoising step away from Gaussian noise. No larger budget hands the projector such a plan,
because there the gated steps begin in the middle of the schedule; and the flow-based models never do,
because their projection window is the last $\lceil 0.5\,\nfe \rceil$ solver steps, which at two steps
is the terminal one alone. Projecting the final step costs what it costs everywhere: the post-processing
variant of the same cell, which projects nothing but the final step, runs at $28$\,ms per control step,
$10$\,ms above the unprojected plan. The remaining $180$--$190$\,ms is the first-step projection,
about $45$\,ms per candidate against $2$--$6$\,ms for a plan that is nearly denoised. The baseline's
projector is a sequential quadratic program solved iteratively, and its cost is set by how far from
feasible and how rough the plan it starts from is: a plan that is still mostly noise violates every
constraint at once and the dynamics rows between every pair of steps, and the solver spends its
iterations there.

Two evaluations on different days give the same time on all five seeds. At equal noise level the two
families cost about the same per solve: at $\nfe=20$ the diffusion baseline projects eleven steps from
the middle of the schedule at about $8.5$\,ms per solve, and instantaneous-velocity matching ten at about
$7$\,ms. Whether a flow-based model's first-step plan --- a point on the straight path from noise to
data --- projects more cheaply than the diffusion model's at the same step is not measured here: at two
steps the flow-based window excludes that step. The measurement that would show it is the same cell with
the projection window opened to both steps on the flow-based side.
The ordering of this benchmark does not rest on this cell: every flow-based row is cheaper than
every diffusion row at equal success with constraint satisfaction, and the operating points are at one
evaluation, where both families project a single, final step.

\subsubsection{Endpoint and Per-Step Projection}\label{sec:res:avoiding:projectors}

The models above were read twice, unprojected and under the per-step projector. This subsection holds the
model fixed and varies the projector, and it is made on matched work: every comparison gives both methods
four candidate plans and reads them on the same training seed, the same three geometries and the same
episodes, so that the two projectors are priced on the same thing. It is read at two budgets. At the
default threshold of $0.5$, $\nfe=3$ is the smallest budget at which endpoint projection has a guiding step
(\autoref{sec:res:constraints:degenerate}); at the threshold of $1.0$ the $\nfe=3$ rows were run at, $\nfe=2$
already has one and $\nfe=3$ two. At $\nfe=10$ both methods have several. On the two average-velocity models the endpoint sampler also
queries the network at a zero interval where the plain sampler queries it at the step width
(\autoref{sec:setup:protocol:guiding}), so on those models a difference between the rows belongs to the
sampler and the point of projection together. \autoref{tab:avoiding-projectors} prints every model at both budgets.
No endpoint-projection evaluation of the diffusion baseline exists on D3IL-avoiding, so no
statement in this subsection extends to it.

\begin{table}[htpb]
  \caption[Endpoint against per-step projection at matched candidates]{Endpoint projection
  \parencite{li2025hardflow} against per-step projection on D3IL-avoiding, both with \textbf{four
  candidate plans}, tightened constraints, the same seed, geometries and episodes on both sides of every
  pair. Both under temporal consistency ($t$) unless marked. At $\nfe=3$ endpoint projection ran at activation threshold $1.0$ (two guiding
  steps), at $\nfe=10$ at $0.5$ (four). S\&C, control steps and the time to compute one action are means
  over three geometries $\times$ two episodes and over the training seeds listed, except CI-MeanFM at
  $\nfe=10$, whose cell holds twenty episodes per geometry. Bold: the better of the pair. Temporal U-Net of
  $4.0$\,M parameters.}
  \label{tab:avoiding-projectors}
  \centering
  \footnotesize
  \setlength{\tabcolsep}{4pt}
  \begin{tabularx}{\linewidth}{@{}L r r L rrr@{}}
    \toprule
    Model & $\nfe$ & Seeds & Projection & S\&C & Steps & ms/step \\
    \midrule
    MeanFM    & 3  & 4 & per-step, $t$ & 0.917 & 59.3 & 148 \\
              &    &       & endpoint      & \textbf{1.000} & 58.9 & \textbf{74.5} \\
    CI-MeanFM, $\alpha_{\mathrm{end}}=0.2$ & 3 & 4 & per-step, $t$ & 0.958 & 60.0 & 146 \\
              &    &       & endpoint      & \textbf{1.000} & 60.0 & \textbf{76.1} \\
    FM        & 3  & 4     & per-step, $t$ & 1.000 & 61.2 & \textbf{64.1} \\
              &    &       & endpoint      & 1.000 & 60.8 & 71.5 \\
    \midrule
    MeanFM    & 10 & 1     & per-step, $t$ & 1.000 & 64.2 & 362 \\
              &    &       & endpoint      & 1.000 & 64.8 & \textbf{176} \\
    CI-MeanFM, $\alpha_{\mathrm{end}}=0.2$ & 10 & 1 & per-step, $t$ & 0.983 & 60.5 & 386 \\
              &    &       & endpoint      & \textbf{1.000} & 60.7 & \textbf{176} \\
    FM        & 10 & 1     & per-step, $c$ & 1.000 & 63.2 & 199 \\
              &    &       & endpoint      & 1.000 & 64.0 & \textbf{167} \\
    \bottomrule
  \end{tabularx}
\end{table}

At $\nfe=3$, on the same four training seeds for every model, endpoint projection reaches $1.000$ success
with constraint satisfaction on all three. On the two average-velocity models it is ahead on the
constraint --- $1.000$ against $0.917$ for analytic average-velocity matching and against $0.958$ for the
consistency-interpolated model --- at level control steps and about half the time per step, $74.5$ against
$148$\,ms and $76.1$ against $146$\,ms, so there it Pareto-dominates per-step projection. The factor of two
is the projector: the predicted endpoint is close to the feasible set at every step, so the single solve
needs few iterations, where the per-step projector solves one program per candidate at each active step.
On instantaneous-velocity matching both projectors reach $1.000$ and per-step projection is the cheaper,
$64.1$ against $71.5$\,ms, with endpoint projection $0.4$ control steps shorter: a trade-off. The per-step
time differs between the models' evaluations ($64$ against $146$--$148$\,ms), so the table is read within
a model.

That result does not move the operating point of this task, because $\nfe=3$ is not on its cost frontier.
\autoref{fig:avoiding-tradeoff} and \autoref{tab:avoiding-dpcc-protocol} place the frontier at one and
two network evaluations: per-step projection at $\nfe=1$ already reaches $0.967$ and $1.000$ for the two
average-velocity models at about $18$\,ms per control step, on five seeds. Five of the six $\nfe=3$ points of
\autoref{tab:avoiding-projectors} are dominated on all three columns by CI-MeanFM at one evaluation, on four
seeds against five; the sixth, MeanFM under endpoint projection, matches that point's success within $0.3$
control steps at four times its time per step. The budget of three therefore offers similar control-step
counts at three to eight times the planning time, and it is dropped from this benchmark's operating point
--- with it the one budget at which endpoint projection would be the better projector at a price that
matters.

At $\nfe=10$, where the comparison is at the same budget on one training seed, endpoint projection is level
or ahead on the constraint and cheaper on all three models. On analytic average-velocity matching both reach
$1.000$, at $176$ against $362$\,ms; on the consistency-interpolated model endpoint projection reaches
$1.000$ at $176$\,ms against $0.983$ at $386$\,ms; on instantaneous-velocity matching both reach $1.000$, at
$167$ against $199$\,ms. Under the other two rules the picture is mixed. Under random selection analytic
average-velocity matching reaches $0.500$ with per-step projection and $1.000$ with endpoint projection, the
largest gap in the table, while the consistency-interpolated model scores $0.833$ and $0.883$ with endpoint
projection against $0.850$ and $0.900$ with per-step projection, level within one episode. The
cumulative-cost rule lengthens the endpoint paths, to $88$ to $111$ control steps, because with four
candidates it can select a stalled plan. The $\nfe=10$ rows are one training seed: six episodes per
rule for analytic average-velocity and instantaneous-velocity matching, sixty for the consistency-interpolated
model, whose cell was evaluated with twenty episodes per geometry. The $\nfe=3$ rows are four training
seeds. Nothing here is read at the protocol of \ac{DPCC}, and nothing here changes a selection made in
\autoref{sec:res:avoiding:dpccprotocol}.

The consequence for this task is that endpoint projection is not needed here, and the reason is the budget
rather than the method. At one network evaluation the two average-velocity models already reach $0.967$
and $1.000$ at about $18$\,ms per control step with the per-step projector, and instantaneous-velocity
matching reaches $1.000$ at one and at two (\autoref{tab:avoiding-dpcc-protocol}); there is no shortfall
left for a second projection method to close. Endpoint projection needs sampling steps to guide, and at
one and two evaluations it has none at the default threshold (\autoref{sec:method:degenerate},
\autoref{sec:res:constraints:degenerate}); on instantaneous-velocity matching a run at $\nfe=2$ that handed
both methods the same candidate plans produced bit-identical rollouts, while on the average-velocity models
the two arms differ even there (\autoref{sec:setup:protocol:guiding}). Where it does act, at three and at ten
evaluations, it is level or ahead on the constraint at the same budget, and cheaper on every model except
instantaneous-velocity matching at three evaluations --- and that is the result carried to the tasks that are
run at those budgets, D3IL-aligning and the quadrotor scenes.

On D3IL-avoiding endpoint projection therefore has no task to win: the per-step projector at one or two
evaluations already reaches $0.967$--$1.000$ at about $18$\,ms, and the budget of three, the first at which
endpoint projection can act at the default threshold, reaches similar control steps at three to eight times
the time, so it lies off the cost frontier of \autoref{sec:res:avoiding:pareto}, which admits success within
$0.05$ of the best. On D3IL-aligning, which needs twenty
evaluations, endpoint projection keeps more contexts violation-free than per-step projection
(\autoref{sec:res:aligning:projection}).

\subsubsection{Steps at Which Projection Guides Sampling}\label{sec:res:constraints:degenerate}

The guiding step count is defined in \autoref{sec:method:degenerate} and the convention for reporting a
configuration that has none in \autoref{sec:setup:protocol:guiding}. This section reports what the
evaluations measured against it.

The count holds. With the activation threshold at $1.0$ the evaluation ran with one guiding step at
$\nfe=2$, two at $\nfe=3$ and four at $\nfe=5$, and the number of programs it actually solved grew as
$1:1.51:2.52$ against the predicted $1:1.5:2.5$, so the configurations reported below are the ones the
prediction says they are.

Two settings in this chapter have no guiding step, and neither supports a statement about endpoint
projection. A budget of $\nfe=1$ has none at any threshold, which is why D3IL-avoiding at its operating
budget is reported with per-step projection only and why UAV-corridor has no endpoint row at that budget;
at the default threshold of $0.5$, $\nfe=2$ has none either, which is where the bit-identical rollouts
of \autoref{sec:res:avoiding:projectors} come from and why \autoref{tab:avoiding-projectors} starts at
$\nfe=3$.

\subsection{Conclusion}\label{sec:res:avoiding:conclusion}

D3IL-avoiding carries the central claim of this thesis, and the results support it: CI-MeanFM
Pareto-dominates the diffusion model of \ac{DPCC} on its own benchmark, at its protocol and with its own
temporal U-Net of $4.0$\,M parameters, and MeanFM matches the saving at one episode in thirty
(\autoref{tab:avoiding-conclusion}). CI-MeanFM at one network evaluation, with temporal consistency and
per-step projection, reaches $1.000$ success with constraint satisfaction as the baseline of record does, in
$59.2$ control steps against $70.1$ and at $18.1$\,ms per step against $553.4$, a thirtieth of the time.
Analytic average-velocity matching at the same budget and projector reaches $0.967$ in $58.6$ steps at
$18.3$\,ms: the same saving, one episode in thirty fewer within the constraints --- a trade-off, not a
dominance.

Instantaneous-velocity matching at one evaluation dominates the baseline as well, at $1.000$ in $67.0$ control
steps and $17.3$\,ms, so much of the saving belongs to a flow-based sampler at a small budget; the
average-velocity models add to it the fewest control steps of the benchmark, about eight fewer than
instantaneous-velocity matching for about a millisecond more per step. The saving does not rest on the
baseline's twenty steps. Retrained at two denoising steps, the baseline also reaches $1.000$, in $61.2$
control steps at $211.4$\,ms, and against that configuration the flow-based models are twelve times faster
rather than thirty. Its two-step time is six times its one-step time for twice the denoising steps, because
of what its projector is handed: at two steps it projects a plan one denoising step away from noise, which costs about $45$\,ms per candidate to repair against $2$ to
$6$\,ms for a nearly denoised plan, while the flow-based models at one or two evaluations project only their
terminal step. What the frontier shows is therefore a region, not a single point: one or two evaluations of
a flow-based model under temporal consistency, inside which the remaining choices trade a control step
against a millisecond. Cumulative projection cost, the baseline's best rule, is not theirs: it occasionally
selects a stalled plan and costs them control steps.

\begin{table}[htpb]
  \caption[Supported combinations on D3IL-avoiding]{The supported model--projection combinations on
  D3IL-avoiding and the baseline of record, at the protocol of \ac{DPCC}. All use tightened constraints
  and the temporal U-Net of $4.0$\,M parameters. The episodes column states the number per geometry;
  every row uses five training seeds and all three geometries.}
  \label{tab:avoiding-conclusion}
  \centering
  \footnotesize
  \begin{tabularx}{\linewidth}{@{}L r l c rrrr@{}}
    \toprule
    Model & $\nfe$ & Projection & Rule & Episodes & S\&C & Steps & ms/step \\
    \midrule
    \mainconfig{\selectedmark\ MeanFM} & \mainconfig{1} & \mainconfig{per-step} &
      \mainconfig{$t$} & \mainconfig{10} & \mainconfig{0.967} & \mainconfig{58.6} & \mainconfig{18.3} \\
    CI-MeanFM & 1 & per-step & $t$ & 10 & 1.000 & 59.2 & 18.1 \\
    \mainconfig{\baselinemark\ Diffusion} & \mainconfig{20} & \mainconfig{per-step} &
      \mainconfig{$c$} & \mainconfig{10} & \mainconfig{1.000} & \mainconfig{70.1} & \mainconfig{553.4} \\
    \bottomrule
  \end{tabularx}
\end{table}

Endpoint projection is not needed at either low-budget operating point. At one evaluation it has no
guiding step, and at two it has none at the default threshold and one at a threshold of $1.0$; raising the
budget so that it can guide sampling gives up the
low-step advantage that decides this benchmark. Where it can act, at three evaluations, it halves the time
per control step of the two average-velocity models at a success with constraint satisfaction of $1.000$,
because the endpoint it projects is already close to the feasible set: the other face of the two-step
baseline's cost, since the price of a projection is set by how far from feasible the plan it is handed is.
Those budgets lie off the cost frontier, so no endpoint-projection configuration improves on the operating
point at one evaluation with per-step projection: on its own benchmark the operating point leads both the
diffusion model of \ac{DPCC} and endpoint projection. D3IL-aligning and the UAV scenes below test endpoint
projection at larger budgets.

The appendix carries more of this benchmark: the executed end-effector paths of the two average-velocity
models at their operated budgets (\autoref{app:avoiding-paths}) and the evaluation at twenty episodes per
seed (\autoref{app:avoiding-twenty}).

\section{Vision-Conditioned Alignment}\label{sec:res:visual}

All results in this section use one training seed and the same ten contexts, and every comparison is made
context by context (\autoref{sec:setup:protocol:aligning}). The box starts on average $0.4530$\,m
from its target. The generative models are compared on the final box-to-target distance of plans executed
without projection on the untightened set, where the constraint set has no bearing --- the plan is the
same plan --- so all four models are compared at once; constraint violations and the time one action costs
are compared in \autoref{sec:res:aligning:projection}, on the tightened set, where the evaluation's
box--obstacle guard freezes one context and the unprojected rows are read within that evaluation.
The ten contexts are drawn from the training contexts of the D3IL task, not from its held-out set,
so these results measure how well each model plans for initial poses of the kind it was trained on.

\subsection{Generative Models}\label{sec:res:aligning:models}

One model moves the box most of the way to its target on this task and the other three barely move it at all
(\autoref{tab:va-models}). The median is the column to read there: the spread over the ten contexts is
of the order of the mean itself, because a context in which the box is not touched contributes the whole
starting distance and one in which it is delivered contributes almost nothing. Analytic average-velocity matching ends a median $0.0741$\,m from the target,
from a starting distance of $0.4530$\,m --- it closes $84\,\%$ of the distance --- and it is the only model
that moves the box in every context. The other three end between $0.39$ and $0.46$\,m. The two
flow-based models close $10$--$14\,\%$ of the distance, while the diffusion median is $2\,\%$ farther
away than the box started; each leaves it untouched in some contexts.

The margins behind that ordering are one-sided. Analytic average-velocity matching ends closer than the
diffusion baseline in all ten contexts, by $0.3744$\,m on average, at $0.64$ of the baseline's time per
control step, and closer than instantaneous-velocity matching in nine of ten, by $0.2333$\,m. Between
instantaneous-velocity matching and the baseline there is no ordering: their medians are $0.4085$ and
$0.4619$\,m, and each ends closer in about half the contexts. Instantaneous-velocity matching leaves the
box unmoved in four of ten contexts and the diffusion baseline in five of ten.

\begin{table}[htpb]
  \caption[Alignment: the generative models without projection]{The generative models on
  D3IL-aligning without projection: the three flow-based models at $\nfe \in \{2, 10, 20, 100\}$ and the
  diffusion baseline at its trained budget, printed last. Final box-to-target distance over the ten
  contexts from a mean initial $0.4530$\,m, in brackets the reduction of the median relative to that mean,
  in per cent (\autoref{sec:setup:metrics:aligning}); \emph{unmoved}: contexts in which the box does not move; mean
  and time as mean $\pm$ sample standard deviation over the contexts.
  Temporal U-Net of $4.0$\,M parameters, FiLM conditioning, one training seed. \selectedmark\ the
  selected configuration; \baselinemark\ the diffusion model of \ac{DPCC}.}
  \label{tab:va-models}
  \centering
  \footnotesize
  \begin{tabularx}{\linewidth}{@{}L r rrr r@{}}
    \toprule
    Model & $\nfe$ & Median [m] & Mean [m] & Unmoved & ms/step \\
    \midrule
    MeanFM     & 2  & 0.2779 (39\,\%) & $0.368 \pm 0.301$ & 0/10  & $28.4 \pm 2.0$ \\
    CI-MeanFM, $\alpha_{\mathrm{end}}=0.2$ & 2  & 0.3738 (17\,\%) & $0.357 \pm 0.201$ & 2/10  & $28.7 \pm 5.7$ \\
    CI-MeanFM, $\alpha_{\mathrm{end}}=0.05$ & 2  & 0.3978 (12\,\%) & $0.336 \pm 0.166$ & 0/10  & $27.6 \pm 2.0$ \\
    FM         & 2 & 0.4257 (6\,\%) & $0.375 \pm 0.133$ & 3/10  & $32.5 \pm 2.5$ \\
    \midrule
    MeanFM     & 10 & 0.1194 (74\,\%) & $0.173 \pm 0.153$ & 1/10 & $99.2 \pm 1.9$ \\
    CI-MeanFM, $\alpha_{\mathrm{end}}=0.2$ & 10 & 0.4247 (6\,\%) & $0.399 \pm 0.094$ & 2/10 & $98.5 \pm 4.0$ \\
    FM         & 10 & 0.4224 (7\,\%) & $0.415 \pm 0.085$ & 4/10 & $145.1 \pm 2.0$ \\
    \midrule
    \mainconfig{\selectedmark\ MeanFM} & \mainconfig{20} & \mainconfig{0.0741 (84\,\%)} & \mainconfig{$0.093 \pm 0.073$} & \mainconfig{0/10} & \mainconfig{$190.5 \pm 2.0$}  \\
    CI-MeanFM, $\alpha_{\mathrm{end}}=0.2$ & 20 & 0.3897 (14\,\%) & $0.325 \pm 0.224$ & 2/10  & $189.1 \pm 5.6$ \\
    CI-MeanFM, $\alpha_{\mathrm{end}}=0.05$ & 20 & 0.4488 (1\,\%) & $0.436 \pm 0.283$ & 0/10  & $187.3 \pm 2.0$ \\
    FM         & 20 & 0.4085 (10\,\%) & $0.327 \pm 0.160$ & 4/10  & $295.8 \pm 1.9$ \\
    \midrule
    MeanFM     & 100 & 0.0673 (85\,\%) & $0.136 \pm 0.148$ & 0/10 & $892.4 \pm 1.8$ \\
    CI-MeanFM, $\alpha_{\mathrm{end}}=0.2$ & 100 & 0.1791 (60\,\%) & $0.327 \pm 0.323$ & 1/10 & $902.2 \pm 1.5$ \\
    FM         & 100 & 0.4247 (6\,\%) & $0.352 \pm 0.180$ & 0/10 & $1425.0 \pm 1.8$ \\
    \midrule
    \mainconfig{\baselinemark\ Diffusion} & \mainconfig{20} & \mainconfig{0.4619 ($-2$\,\%)} & \mainconfig{$0.468 \pm 0.399$} & \mainconfig{5/10} & \mainconfig{$298.3 \pm 2.0$} \\
    \bottomrule
  \end{tabularx}
\end{table}

The ordering rests on analytic average-velocity matching's margins, which hold in nine and ten of
ten contexts; the level result for instantaneous-velocity matching and diffusion is the absence of a
difference at ten contexts, not evidence that they are equal.

\subsubsection{Step Budget and the Cost Frontier}\label{sec:res:aligning:budget}

Analytic average-velocity matching improves with the budget: its median final distance falls from
$0.2779$\,m at $\nfe=2$ to $0.1194$\,m at $\nfe=10$, $0.0741$\,m at $\nfe=20$, and
$0.0673$\,m at $\nfe=100$ (\autoref{tab:va-models}). The last fivefold increase in evaluations reduces the median distance by
$0.0068$\,m while raising the time per control step from $190.5$ to $892.4$\,ms. Unlike D3IL-avoiding,
where one network evaluation already saturates the task (\autoref{sec:res:state}), alignment rewards
additional evaluations.

The three flow-based models are measured at $\nfe=100$ as well, and the ordering among them is the same as
at twenty: analytic average-velocity matching first ($0.0741$ and $0.0673$\,m), then the
consistency-interpolated model ($0.3897$ and $0.1791$\,m), with instantaneous-velocity matching last
($0.4085$ and $0.4247$\,m). A hundred evaluations more than halve the consistency-interpolated model's
median and leave instantaneous-velocity matching where it was. On the complete ladder
instantaneous-velocity matching closes $6$ to $10\,\%$ of the distance at every budget from two to a hundred,
and the consistency-interpolated model $17$, $6$, $14$ and $60\,\%$, so at ten evaluations it is at the
level of instantaneous-velocity matching. Analytic average-velocity matching leads at every budget: at ten
evaluations it ends closer than instantaneous-velocity matching at less time per control step ($99.2$
against $145.1$\,ms), and closer than the consistency-interpolated model at the same time ($98.5$\,ms). The diffusion baseline is compared at the
budget it was trained with. The ladder starts at two evaluations: the analytic model closes $39\,\%$ of
the distance there against $84\,\%$ at twenty, so a smaller budget is not an operating point on this
task.

Putting the outcome and its price on the same axes gives this task its cost frontier
(\autoref{fig:aligning-tradeoff}). It is not the frontier of \autoref{fig:avoiding-tradeoff} and cannot
be: on D3IL-avoiding every eligible configuration solves the task and only the two costs differ, whereas
here nothing is solved or failed --- what a model delivers is a distance, and the control-step count is
set by the episode limit rather than by arrival (\autoref{tab:eval}). The frontier is therefore the
outcome itself against the time one action costs, and a configuration is on it when nothing else both
ends closer to the target and computes faster.

Three marks of the drawing carry a reading rule. The outcome axis is the reduction of the median final
distance relative to the mean initial distance, the bracketed number of \autoref{tab:va-models}, so the
dashed line at $0\,\%$ is a median final distance equal to the mean initial one and $100\,\%$ at the top is
a median box delivered; two points
that would print on top of each other are displaced by a few pixels, which moves the drawing and not the
data; and the boxed arrow is the direction in which both axes improve.

Four of the thirteen configurations are on it, and all four are analytic average-velocity matching, at
$\nfe=2$, $10$, $20$ and $100$. Every other model at every budget is dominated, and the way they are dominated
is the result: at $\nfe=20$ the consistency-interpolated objective, instantaneous-velocity matching and
the diffusion baseline all sit close to the $0\,\%$ line, at between $189$
and $298$\,ms per control step, while the analytic objective at $\nfe=2$ --- a seventh to a tenth of
that cost --- has already closed two fifths of the distance. MeanFM's own $\nfe=100$ point reduces its median distance
by only $0.0068$\,m at $4.7$ times its $\nfe=20$ cost.

\begin{figure}[htpb]
  \centering
  \includegraphics[width=0.92\linewidth]{fig_aligning_tradeoff}
  \caption[Alignment: outcome against cost]{D3IL-aligning without projection: the reduction of the
  median final box-to-target distance relative to the mean initial distance, in per cent
  (\autoref{sec:setup:metrics:aligning}), against the time to compute one action on a log axis; every cell
  of \autoref{tab:va-models} except the two $\alpha_{\mathrm{end}}=0.05$ rows, one training seed. The
  dashed red line is the $0\,\%$ reference. Rings mark the non-dominated configurations, joined by the
  dashed staircase; beside each point is its step budget; the boxed arrow points where both axes improve.
  The outcome axis is logarithmic in the median distance left to the target; the tick labels are the
  reduction.}
  \label{fig:aligning-tradeoff}
\end{figure}

\subsubsection{Consistency Interpolation across the Step Budget}\label{sec:res:aligning:consistency}

The consistency-interpolated objective behaves differently from the analytic one on this task: it does
not improve when it is given a larger budget. At two network evaluations the two are level --- each ends
closer in five of the ten contexts. At twenty, where the analytic objective has moved the box most of the
way to its target, the consistency-interpolated one has not moved at all from where it was at two
evaluations, and it ends further away in seven of ten contexts. Raising the budget spends time and buys
nothing, and the same holds for the more aggressive setting $\alpha_{\mathrm{end}}=0.05$
(\autoref{tab:va-models}).

The numbers behind that are flat where the analytic objective's fall. With
$\alpha_{\mathrm{end}}=0.2$ the median final distance moves from $0.3738$ to $0.3897$\,m between two and
twenty evaluations, and with $\alpha_{\mathrm{end}}=0.05$ from $0.3978$ to $0.4488$\,m, against the analytic
objective's $0.2779$ to $0.0741$\,m; the time per control step is the same for the three at each budget.

The mechanism is in the training target. It is built from a stop-gradient copy of the network itself, so
part of the network's own error enters the field it learns, and that part does not shrink when the solver
takes smaller steps (\autoref{sec:method:alphaflow}). A task that rewards more evaluations therefore
exposes it, as this one does; a task already saturated at one evaluation, as D3IL-avoiding is, does not.
The level result at $\nfe=2$ is in distance only. Counting the contexts in which any
configuration reaches the target, analytic average-velocity matching succeeds in seven and
consistency-interpolated average-velocity matching in two.

\subsection{Projection Methods}\label{sec:res:aligning:projection}

D3IL-aligning is the task on which the generative model needs many network evaluations: analytic
average-velocity matching comes closest to the target, and $\nfe=20$ is its operating budget --- the
$\nfe=100$ point ends only $0.0068$\,m closer at $4.7$ times the cost (\autoref{sec:res:aligning:budget})
--- so that is the model and the budget the projection is put on. The projected results below are all read
on it.

\paragraph{The activation threshold is compressed, because the budget is high.}
Everything about the threshold on this task follows from one fact: D3IL-aligning is the task that needs a
large step budget, and the default of \ac{DPCC}, $\eta=0.5$, was set for small ones. The projector runs
on the last $\lceil \eta\nfe \rceil$ sampling steps, so at a fixed share $\eta$ its cost grows with the
budget: every extra sampling step in the active window is another program solved, for each of the four
candidate plans, at each of up to four hundred control steps. On D3IL-avoiding at one or two evaluations
the default costs one solve; here, at $\nfe=20$, it would be ten solves per candidate per control step,
and at $\nfe=100$ fifty.

What a solve costs is measured on the two budgets the projector has been run at
(\autoref{tab:va-threshold}). At $\nfe=20$ with $\eta=0.2$ the projector runs on four sampling steps,
sixteen solves per control step, and the projected control step costs $93$\,ms more than the unprojected
one: about $6$\,ms per solve. At $\nfe=10$ with $\eta=0.4$, the same sixteen solves cost $428$\,ms,
$27$\,ms per solve, the plans handed to the projector there being further from their final shape. With
those two figures the default threshold at $\nfe=100$ prices itself: fifty active steps are two hundred
solves per control step, $1.2$--$5.4$\,s on top of the $0.9$\,s the hundred evaluations take, and an
alignment cell is four hundred control steps on ten contexts --- between two and seven hours of computation
for one cell as a lower bound, since the early sampling steps of a hundred-step schedule are noisier than
any step of a twenty-step one. That pair is not run. The compressed threshold at the same budget,
$\eta=0.1$, is forty solves per control step, which the same figures price at $0.2$--$1.1$\,s on top of the
generation; measured, the forty solves add $187$\,ms to the $925$\,ms of the hundred evaluations, under
$5$\,ms per solve (\autoref{tab:va-threshold}).

The threshold is therefore compressed as the budget grows, so that the projector keeps a roughly constant
\emph{number} of guiding steps instead of a constant share: $\eta=0.2$ at $\nfe=20$ and $\eta=0.4$ at
$\nfe=10$, which is three guiding steps at both (\autoref{sec:setup:protocol:guiding}), where on
D3IL-avoiding at one evaluation there were none. Every projected alignment result of this chapter is under
that compressed threshold, and \autoref{tab:va-threshold} is the ladder it was chosen on. What a lower threshold gives up is the early part of sampling, where the sample is still
far from its final shape and the projection has least to hold on to; what it buys is a projected
configuration whose time per action is of the same order as the generation it is attached to.

Each row of \autoref{tab:va-threshold} is the analytic average-velocity model under random selection on
the tightened set, the configuration of the $\nfe=100$ comparison, under \emph{both} projectors: the
per-step projector of \ac{DPCC}, which is what the threshold governs, and endpoint projection at the same
threshold, whose guiding steps are the same count minus the terminal one. The number of guiding steps is
the quantity held roughly constant. The $\nfe=2$, $10$ and $20$ rows are the random-rule cells of
\autoref{tab:va-projection}, the compressed $\nfe=100$ row is its own evaluation, and the uncompressed
$\nfe=100$ row is not run, for the cost set out above.

\begin{table}[htpb]
  \caption[Alignment: the activation threshold across the budget ladder]{The activation threshold
  $\eta$ on D3IL-aligning: MeanFM under random selection, tightened constraints, ten contexts, one
  training seed, under both projectors. \emph{Guiding}: sampling steps the projection acts on minus the
  terminal one, $\lceil \eta\nfe \rceil - 1$ (\autoref{sec:setup:protocol:guiding}). \emph{ms/step}:
  mean over the contexts, unprojected and under each projector;
  \emph{median}: median final box-to-target distance from $0.4530$\,m; \emph{violation-free}: contexts free of
  violations. The $\nfe=2$, $10$ and $20$ rows are the random-rule cells of \autoref{tab:va-projection}.}
  \label{tab:va-threshold}
  \centering
  \footnotesize
  \setlength{\tabcolsep}{3pt}
  \begin{tabularx}{\linewidth}{@{}r r r r rrr rrr@{}}
    \toprule
    & & & & \multicolumn{3}{c}{per-step projection} & \multicolumn{3}{c}{endpoint projection} \\
    \cmidrule(lr){5-7}\cmidrule(l){8-10}
    $\nfe$ & $\eta$ & Guiding & \shortstack{ms/step\\unproj.} & ms/step & Median [m] & \shortstack{Violation-\\free} & ms/step & Median [m] & \shortstack{Violation-\\free} \\
    \midrule
    2   & 0.5 & 0  & $24.6$  & $46.2$  & 0.275 & 8/10  & \multicolumn{3}{l}{--- no guiding step} \\
    10  & 0.4 & 3  & $89.7$  & $517.5$ & 0.434 & 7/10  & $192.5$ & 0.242 & 10/10 \\
    20  & 0.2 & 3  & $172.5$ & $265.9$ & 0.179 & 9/10  & $275.3$ & 0.197 & 10/10 \\
    100 & 0.1 & 9  & $924.6$ & $1112.0$ & 0.248 & 8/10 & $1107.0$ & 0.153 & 8/10 \\
    100 & 0.5 & 49 & \multicolumn{7}{l}{\emph{not run}: two hundred solves per control step (text)} \\
    \bottomrule
  \end{tabularx}
\end{table}

Along the ladder the compressed threshold holds the projected control step at $1.9$, $5.8$, $1.5$ and
$1.2$ times the unprojected one at $\nfe=2$, $10$, $20$ and $100$ under per-step projection, and at $2.1$,
$1.6$ and $1.2$ times under endpoint projection where it has guiding steps: at a budget of a hundred,
nine guiding steps add a fifth to the time of the generation. Endpoint projection keeps ten of ten contexts
violation-free at budgets of ten and twenty and eight at a hundred; per-step projection keeps seven to nine at
every budget. Under either projector the median ends further from the target than the unprojected plan at
budgets of ten and twenty: $0.179$ and $0.197$\,m against $0.090$\,m at $\nfe=20$, $0.434$ and $0.242$\,m
against $0.142$\,m at $\nfe=10$. At a budget of two the per-step median, $0.275$\,m, is closer than the
unprojected $0.320$\,m, and at a hundred endpoint projection ends at $0.153$\,m, the smallest median of any
projected row of the ladder.

\autoref{tab:va-projection} compares the two projection methods on the analytic average-velocity model,
on tightened constraints, under all three selection rules.

Every cell of that table, and of every other alignment table in this chapter, is the same ten contexts.
Some of the evaluations behind them were run over thirty contexts rather than ten; those ten are a subset
of those thirty in every case, so each cell is restricted to them. The section is then paired on one
context set throughout, as \autoref{sec:setup:protocol:aligning} says it is, rather than setting a
ten-context result beside a thirty-context one.

The $\nfe=2$ block has per-step rows only. At its activation threshold of $0.5$ a two-step budget has no
guiding step --- $\lceil \eta\nfe\rceil - 1 = 0$ (\autoref{sec:setup:protocol:guiding},
\autoref{sec:res:constraints:degenerate}) --- so endpoint projection does not exist there as a method: what
the evaluation ran under that name at $\nfe=2$ is projection after sampling, and it is not printed as an
endpoint row. The smallest budget at which the two methods can be compared is $\nfe=3$.

\begin{table}[htpb]
  \caption[Alignment: the two projection methods]{D3IL-aligning, MeanFM with each
  projection method, tightened constraints, ten contexts, one training seed. \emph{Violation-free}:
  contexts with no violating control step; \emph{violating steps}: per context; \emph{moved}: contexts
  in which the box moves; \emph{median} and \emph{mean}: final box-to-target distance over the ten
  contexts, the median with its reduction relative to the mean initial $0.4530$\,m in brackets, the mean $\pm$
  sample standard deviation; ms/step likewise. Rules $r$, $c$, $t$ as in
  \autoref{tab:avoiding-dpcc-protocol}. \selectedmark\ the selected configuration.}
  \label{tab:va-projection}
  \centering
  \footnotesize
  \setlength{\tabcolsep}{4pt}
  \begin{tabularx}{\linewidth}{@{}r L c rrrrrr@{}}
    \toprule
    $\nfe$ & Projection & Rule & \shortstack{Violation-\\free} & \shortstack{Violating\\steps} & Moved & Median [m] & Mean [m] & ms/step \\
    \midrule
    20 & none      & ---  & 2/10  & 31.6 & 9/10 & 0.090 (80\,\%) & $0.143 \pm 0.134$ & $172.5 \pm 60.6$ \\
       & per-step  & $r$  & 9/10  & 2.7  & 8/10 & 0.179 (61\,\%) & $0.200 \pm 0.150$ & $265.9 \pm 101.9$ \\
       &           & $c$  & 9/10  & 4.3  & 6/10 & 0.333 (27\,\%) & $0.306 \pm 0.164$ & $325.3 \pm 165.8$ \\
       &           & $t$  & 8/10  & 1.5  & 7/10 & 0.228 (50\,\%) & $0.240 \pm 0.165$ & $323.2 \pm 175.0$ \\
       & \mainconfig{endpoint}  & \mainconfig{\selectedmark\ $r$} & \mainconfig{10/10} & \mainconfig{0.0} & \mainconfig{8/10} & \mainconfig{0.197 (57\,\%)} & \mainconfig{$0.211 \pm 0.144$} & \mainconfig{$275.3 \pm 102.4$}  \\
       &           & $c$  & 10/10 & 0.0  & 6/10 & 0.323 (29\,\%) & $0.286 \pm 0.190$ & $282.3 \pm 104.1$ \\
       &           & $t$  & 9/10  & 0.6  & 7/10 & 0.230 (49\,\%) & $0.255 \pm 0.154$ & $301.5 \pm 128.6$ \\
    \midrule
    10 & none      & ---  & 1/10  & 35.5 & 8/10 & 0.142 (69\,\%) & $0.199 \pm 0.179$ & $89.7 \pm 31.5$ \\
       & per-step  & $r$  & 7/10  & 12.3 & 8/10 & 0.434 (4\,\%) & $0.420 \pm 0.352$ & $517.5 \pm 344.3$ \\
       &           & $c$  & 6/10  & 21.3 & 6/10 & 0.411 (9\,\%) & $0.352 \pm 0.174$ & $524.3 \pm 417.2$ \\
       &           & $t$  & 7/10  & 5.0  & 4/10 & 0.408 (10\,\%) & $0.271 \pm 0.217$ & $433.5 \pm 224.6$ \\
       & endpoint  & $r$  & 10/10 & 0.0  & 8/10 & 0.242 (47\,\%) & $0.236 \pm 0.126$ & $192.5 \pm 72.3$ \\
       &           & $c$  & 8/10  & 0.6  & 4/10 & 0.425 (6\,\%) & $0.307 \pm 0.198$ & $204.2 \pm 79.4$ \\
       &           & $t$  & 8/10  & 2.3  & 6/10 & 0.255 (44\,\%) & $0.284 \pm 0.161$ & $222.1 \pm 99.0$ \\
    \midrule
    2  & none      & ---  & 5/10  & 46.1 & 8/10 & 0.320 (29\,\%) & $0.323 \pm 0.279$ & $24.6 \pm 8.6$ \\
       & per-step  & $r$  & 8/10  & 11.6 & 9/10 & 0.275 (39\,\%) & $0.286 \pm 0.210$ & $46.2 \pm 18.4$ \\
       &           & $c$  & 8/10  & 6.5  & 5/10 & 0.434 (4\,\%) & $0.329 \pm 0.186$ & $47.3 \pm 19.1$ \\
       &           & $t$  & 10/10 & 0.0  & 7/10 & 0.221 (51\,\%) & $0.250 \pm 0.177$ & $43.0 \pm 16.0$ \\
       & endpoint  & \multicolumn{7}{l}{--- none: no guiding step at $\nfe=2$ (\autoref{sec:setup:protocol:guiding})} \\
    \bottomrule
  \end{tabularx}
\end{table}

At $\nfe=20$, endpoint projection keeps all ten contexts free of violations with random selection and
with cumulative projection cost, and nine with temporal consistency; per-step projection keeps nine, nine
and eight. Together with endpoint projection and random selection at $\nfe=10$, these are the only
configurations at ten or more evaluations that keep every context free of violations while the box is
still moved.
The time per control step is the same for both methods at this budget, $275$ against $266$\,ms with random
selection: what endpoint projection adds at $\nfe=20$ is the removal of the last violations.

At $\nfe=10$ the difference is larger. Endpoint projection with random selection again keeps all ten
contexts free of violations, where per-step projection keeps seven, and it takes $193$ against $518$\,ms
per control step, less in every context whose time differs. With the other two rules it is faster by
$320$ and $211$\,ms per control step and keeps at least as many contexts free of violations.

Of the two budgets, $\nfe=20$ is the operating point, because the generative model itself is closest to
the target there; $\nfe=10$ reaches the same constraint satisfaction with endpoint projection in
$70\,\%$ of the time.
Repeated evaluations of configurations with projection vary by about $0.4$\,m in final
distance, so the median distances of the projected rows are not ranked against each other.

The two-evaluation block says one thing: the budget at which this task can be run cheaply is not the
budget at which the generative model comes closest to solving it. Per-step projection under temporal consistency keeps all
ten contexts free of violations there at $43$\,ms per control step, which nothing at the two larger
budgets does at that price, but the median ends $0.22$--$0.43$\,m from the target, against $0.09$\,m
unprojected at $\nfe=20$.

\subsubsection{Endpoint Projection across the Models}\label{sec:res:aligning:projection:models}

The comparison above is one generative model. \autoref{tab:va-projection-models} runs it across the
models at the operating point this section has settled on --- $\nfe=20$ and the compressed threshold
$\eta=0.2$ --- under every selection rule, and it holds in the same direction. The budget ladder below
twenty was the subject of \autoref{tab:va-projection} and \autoref{tab:va-threshold}; it is not repeated
here, because the model that solves the task does so at twenty and the threshold study has already priced
the other budgets.

All three flow-based models carry both projectors at $\nfe=20$ on the tightened set, and the diffusion
baseline carries per-step projection; endpoint projection is not defined for it. The comparison before
projection in \autoref{tab:va-models} and the one after projection here therefore rest on the same four
models.
Nine comparisons run the two methods on the same model, the same budget and the same rule. Endpoint
projection keeps more contexts free of violations in seven of them and per-step projection in two,
instantaneous-velocity matching and the consistency-interpolated model under temporal consistency, each by
one context. A second evaluation of one projected cell --- analytic average-velocity matching at
$\nfe=2$, per-step, random selection --- moved it by one context and $3$\,cm in median, so a one-context
difference is within what a re-run changes.

That is the reverse of D3IL-avoiding, and the reason is the budget rather than the task. This task is
run at ten and twenty network evaluations because the generative model needs them
(\autoref{sec:res:aligning:budget}), so the endpoint solve has several guiding steps to act on, where on
D3IL-avoiding at one evaluation it had none (\autoref{sec:res:avoiding:projection}). Given steps to act
on, the method does what it is for: it removes violations from the plan before the plan is finished,
rather than correcting each step after the fact.

What it does not do is dominate. Its advantage is on the constraint and not on the task: across the same
nine comparisons it moves the box in more contexts than per-step projection in one, in as many in four,
and in fewer in four, and its median final distances are level. Nor is it cheaper at this budget as a
rule: on the analytic model it costs $0.87$ to $1.04$ of per-step projection's time per control step, on
the consistency-interpolated model $0.71$ to $1.08$, on instantaneous-velocity matching $1.13$ to $1.24$.
Endpoint projection is therefore the better projector on the constraint on this task, at the budget the
task is run at, by a margin that is consistent rather than large, and at a comparable price.

\begin{table}[htpb]
  \caption[Alignment: the two projection methods across the models]{D3IL-aligning, both projection
  methods on the three flow-based models and the diffusion baseline at $\nfe=20$, $\eta=0.2$, tightened
  constraints, the same ten contexts in every row, one training seed. Columns as in
  \autoref{tab:va-projection}; the MeanFM rows repeat the $\nfe=20$ block of that table, and a negative
  share is a box that ends farther from the target than it started. Endpoint projection adds its control to
  the velocity field of a flow-based sampler (\autoref{sec:method:hardflow}); the diffusion model has none,
  so its endpoint row is empty. \selectedmark\ the selected configuration; \baselinemark\ the diffusion model of
  \ac{DPCC}.}
  \label{tab:va-projection-models}
  \centering
  \footnotesize
  \setlength{\tabcolsep}{4pt}
  \begin{tabularx}{\linewidth}{@{}L r c rrrrrr@{}}
    \toprule
    Model & $\nfe$ & Rule & \shortstack{Violation-\\free} & \shortstack{Violating\\steps} & Moved & Median [m] & Mean [m] & ms/step \\
    \midrule
    \multicolumn{9}{@{}l}{\emph{per-step projection}}\\
    MeanFM    & 20 & $r$ & 9/10  & 2.7  & 8/10 & 0.179 (61\,\%) & $0.200 \pm 0.150$ & $265.9 \pm 101.9$ \\
              &    & $c$ & 9/10  & 4.3  & 6/10 & 0.333 (27\,\%) & $0.306 \pm 0.164$ & $325.3 \pm 165.8$ \\
              &    & $t$ & 8/10  & 1.5  & 7/10 & 0.228 (50\,\%) & $0.240 \pm 0.165$ & $323.2 \pm 175.0$ \\
    CI-MeanFM & 20 & $r$ & 4/10  & 44.6 & 8/10 & 0.429 (5\,\%) & $0.432 \pm 0.182$  & $415.8 \pm 268.0$ \\
              &    & $c$ & 5/10  & 10.7 & 5/10 & 0.443 (2\,\%) & $0.388 \pm 0.131$  & $311.4 \pm 130.1$ \\
              &    & $t$ & 6/10  & 23.2 & 7/10 & 0.422 (7\,\%) & $0.353 \pm 0.153$  & $320.5 \pm 153.3$ \\
    FM        & 20 & $r$ & 5/10  & 14.3 & 3/10 & 0.425 (6\,\%) & $0.409 \pm 0.112$  & $409.7 \pm 147.9$ \\
              &    & $c$ & 4/10  & 14.9 & 2/10 & 0.425 (6\,\%) & $0.409 \pm 0.111$  & $406.8 \pm 149.5$ \\
              &    & $t$ & 6/10  & 12.8 & 3/10 & 0.425 (6\,\%) & $0.411 \pm 0.111$  & $402.0 \pm 146.4$ \\
    \mainconfig{\baselinemark\ Diffusion} & \mainconfig{20} & $r$ & 4/10 & 14.8 & 3/10 & 0.462 ($-2$\,\%) & $0.397 \pm 0.160$ & $809.2 \pm 756.0$ \\
              &    & $c$ & 1/10  & 35.0 & 0/10 & 0.456 ($-1$\,\%) & $0.453 \pm 0.037$ & $1069.7 \pm 848.6$ \\
              &    & $t$ & 1/10  & 41.1 & 1/10 & 0.462 ($-2$\,\%) & $0.456 \pm 0.039$ & $766.0 \pm 284.5$ \\
    \midrule
    \multicolumn{9}{@{}l}{\emph{endpoint projection}}\\
    \mainconfig{\selectedmark\ MeanFM} & \mainconfig{20} & \mainconfig{$r$} & \mainconfig{10/10} & \mainconfig{0.0} & \mainconfig{8/10} & \mainconfig{0.197 (57\,\%)} & \mainconfig{$0.211 \pm 0.144$} & \mainconfig{$275.3 \pm 102.4$}  \\
              &    & $c$ & 10/10 & 0.0  & 6/10 & 0.323 (29\,\%) & $0.286 \pm 0.190$ & $282.3 \pm 104.1$ \\
              &    & $t$ & 9/10  & 0.6  & 7/10 & 0.230 (49\,\%) & $0.255 \pm 0.154$ & $301.5 \pm 128.6$ \\
    CI-MeanFM & 20 & $r$ & 8/10  & 2.7  & 7/10 & 0.425 (6\,\%)  & $0.333 \pm 0.183$  & $294.2 \pm 113.3$ \\
              &    & $c$ & 6/10  & 10.1 & 4/10 & 0.456 ($-1$\,\%) & $0.440 \pm 0.076$  & $314.8 \pm 119.6$ \\
              &    & $t$ & 5/10  & 8.6  & 6/10 & 0.434 (4\,\%)  & $0.368 \pm 0.124$  & $344.8 \pm 164.6$ \\
    FM        & 20 & $r$ & 6/10  & 12.0 & 3/10 & 0.443 (2\,\%) & $0.398 \pm 0.132$  & $498.0 \pm 210.4$ \\
              &    & $c$ & 5/10  & 10.4 & 3/10 & 0.425 (6\,\%) & $0.407 \pm 0.109$  & $502.4 \pm 212.0$ \\
              &    & $t$ & 5/10  & 7.5  & 2/10 & 0.425 (6\,\%) & $0.410 \pm 0.110$  & $452.4 \pm 193.4$ \\
    \mainconfig{\baselinemark\ Diffusion} & \mainconfig{20} & \multicolumn{7}{l}{\mainconfig{--- (not defined: no velocity field)}} \\
    \bottomrule
  \end{tabularx}
\end{table}

On the tightened set one of the ten contexts is never attempted: the tightening enlarges the
obstacle until it overlaps the box of context~7 (position~8 in the evaluation's order), the evaluation's
box--obstacle guard holds the box, and
the context scores clean and unmoved in every tightened row. Every $x/10$ of the tightened alignment tables
(\autoref{tab:va-threshold}, \autoref{tab:va-projection}, \autoref{tab:va-projection-models}) is therefore
$x-1$ of nine attempted contexts, and the diffusion baseline's single clean context under $c$ and $t$ is
that one. Comparisons between rows are unaffected.
Consistency-interpolated average-velocity matching and instantaneous-velocity matching barely
move the box on this task at $\nfe=20$ --- medians $14$ and $10\,\%$ of the distance closed, the box unmoved
in two and four of ten contexts (\autoref{tab:va-models}) --- so their rows compare the two
projectors on a plan that does not solve the task. That is what makes them a clean read on the
projector --- the constraint columns are not confounded by how far the box travelled --- and it is also
why the median-distance column of those rows carries no information: every cell sits within a few
centimetres of the starting $0.4530$\,m. The analytic average-velocity rows are the opposite case and are
the ones to read for the median.

The diffusion baseline under per-step projection keeps four contexts free of violations with random
selection and one with each of the other two rules, moves the box in at most three contexts, and ends with
its median slightly farther from the target than the box started. Analytic average-velocity matching under
the same rule Pareto-dominates its best row: nine violation-free contexts against four, $0.179$ against $0.462$\,m,
$266$ against $809$\,ms per control step. The baseline's violations change family under projection:
halfspace and obstacle violations fall to almost none, while executed steps beyond the action bound rise
from $64$ without projection to between $145$ and $408$ over the ten contexts.

The medians of \autoref{tab:va-projection} also hide which contexts changed. The per-context records
show three differences that the medians cannot show.

The constraint column turns over almost completely: eight of the ten contexts carry a violating step
without a projector, one still does under per-step projection, and none does under endpoint projection.
The distances do not turn over with it --- they move around, and mostly in the wrong direction.

The cost of that turnover falls on the contexts the unprojected arm had already delivered. Context~1 ends
$0.065$\,m from the target unprojected and $0.21$ or $0.25$\,m with a projector; context~8 goes from
$0.099$ to $0.28$\,m; context~9 is pushed to within $0.029$\,m unprojected and is then not moved at all.
What the projector pays for constraint satisfaction, it pays in the contexts that were nearly solved.

The exchange is not one-sided. Context~10 improves from $0.143$ to $0.022$\,m and context~4 from
$0.262$ to $0.147$\,m under the same projectors, so what the median reports as a rise from $0.090$ to
$0.197$\,m is not a uniform degradation but a redistribution --- three contexts much worse, two clearly
better, and one, context~7, that no configuration on this task moves at all.

The projected configurations of \autoref{tab:va-projection-models} are compared on the two axes of
\autoref{fig:aligning-tradeoff} in \autoref{fig:aligning-projected-tradeoff}, and what separates the two
drawings is said on the page. Every point of the second is a plan \emph{after} projection on the tightened
set, so it is dearer per control step and, for the same model and budget, further from the target than
the unprojected plan. The faint grey rings are the unprojected cells of analytic average-velocity matching
from \autoref{fig:aligning-tradeoff}, evaluated on the untightened set; the matched reference is the
tightened evaluation's own unprojected row (\autoref{tab:va-projection}), and against it the price reads
from a reduction of $80\,\%$ at $172.5$\,ms to $61\,\%$ at $266$\,ms (from $84\,\%$ at $190$\,ms against the
ring). What that price buys is the constraint: a filled point keeps at least nine of ten
contexts free of violations, which is this task's counterpart of success with constraint satisfaction,
where the unprojected plan keeps one or two. That is the trade this task poses: pursuing no violation
costs completion --- the projected plan ends further from the target --- and allowing some violation
buys completion, the unprojected plan ending closest of all while crossing the keep-out region in seven
contexts of ten and the halfspace in one. The figure places the two ends of that trade on the same axes. Only filled points enter the cost frontier; the hollow ones
remain visible, and on this task they are every point of the two models that barely move the box. The
plotted runs use random selection at $\nfe=20$, the rows of \autoref{tab:va-projection-models}. The
diffusion baseline under per-step projection and the consistency-interpolated model under endpoint
projection keep at most eight contexts free of violations (\autoref{tab:va-projection-models}), below the
nine a filled point needs, so both are hollow and the frontier does not include them.

\begin{figure}[htpb]
  \centering
  \includegraphics[width=0.92\linewidth]{fig_aligning_projected_tradeoff}
  \caption[Alignment after projection: outcome against cost]{D3IL-aligning after projection, on the
  axes of \autoref{fig:aligning-tradeoff}: the cells of \autoref{tab:va-projection-models} under random
  selection, $\nfe=20$, $\eta=0.2$, tightened constraints, one training seed. Colour: MeanFM, CI-MeanFM,
  FM, diffusion; circles per-step projection, squares endpoint projection; the faint dashed grey ring is the
  unprojected MeanFM cell at the same budget from \autoref{fig:aligning-tradeoff}, untightened set. Filled
  points keep at least nine contexts violation-free, hollow ones do not; rings and the dashed staircase mark
  the non-dominated filled points; the dashed line is the $0\,\%$ reference.}
  \label{fig:aligning-projected-tradeoff}
\end{figure}

\subsection{Conclusion}\label{sec:res:aligning:conclusion}

D3IL-aligning repeats the result of D3IL-avoiding on a harder task and adds endpoint projection to it. On
camera observations and an alignment through contact, analytic average-velocity matching again dominates the
diffusion baseline: under per-step projection with random selection it keeps nine contexts free of
violations against the four of the baseline's best projected row, its median box ends $0.179$ against
$0.462$\,m from the target, and it takes $266$ against $809$\,ms per control step. Without projection every
non-dominated configuration of the task is analytic average-velocity matching: at twenty evaluations it closes
$84\,\%$ of the distance, against $14$, $10$ and $-2\,\%$ for the consistency-interpolated model,
instantaneous-velocity matching and the baseline, and at two evaluations, at a seventh to a tenth of their
cost, it has already closed two fifths. The consistency-interpolated model does not improve with the budget
and reaches $60\,\%$ only at a hundred evaluations: its training target is built from a stop-gradient copy of
the network, so part of its own error enters the field it learns, and smaller steps do not remove it. A task
that rewards more evaluations exposes this; D3IL-avoiding, saturated at one evaluation, does not.

At the large budget this task needs, endpoint projection is the better projector on the constraint. With
random selection it keeps all ten contexts free of violations and still moves the box towards its target in
eight, at the time per control step of per-step projection at $\nfe=20$ and at $37\,\%$ of it at $\nfe=10$, and
it keeps more contexts free of violations than per-step projection in seven of nine comparisons. That is the reverse of D3IL-avoiding,
and the reason is the budget rather than the task: at ten and twenty evaluations the endpoint solve has steps
to guide, and it removes violations before the plan is finished. Its advantage is on the constraint and not on
the task, at a comparable price, and the task poses its own trade: pursuing no violation costs completion, the
projected plan ending further from the target than the unprojected one, and the budget at which the task runs
cheaply, two evaluations at $43$\,ms, is not the budget at which the model comes closest to it. No context ends
in position, within the $1.8$\,cm of the task's own test: with endpoint projection the box stops $2.2$ and
$4.7$\,cm from its target in the best two contexts and $13$ to $27$\,cm away in six more, and the unprojected plan's
best two end $2.8$ and $2.9$\,cm away. The conclusion is limited to training-set contexts. The executed end-effector paths of the three projection settings, context by
context, are drawn in \autoref{app:aligning-paths}.

\section{Quadrotor Benchmark}\label{sec:res:uav}

The section is one part per scene, and each scene is read like the two manipulation benchmarks: the
generative models first on the plans they generate, executed without projection; the same cells then under
the projector of \ac{DPCC}, where the cost frontier and the flown paths are; and the two projection methods
against each other last. The scenes come in the order of what they add. UAV-pillars flies the plans of the avoiding planner itself
--- the obstacle field of D3IL-avoiding scaled into an arena of pillars, the same test-time constraints, the
quadrotor as the plant --- so it asks whether a plan the projector repairs on D3IL-avoiding survives execution
by a vehicle. UAV-corridor is the scene whose test-time constraints leave the horizontal plane --- a slide leaned into
the vertical, and a roof --- read one at a time, as the geometries of D3IL-avoiding are, and neither of
which the demonstrations satisfy.
UAV-s-curve is the scene of the two sharp turns. It is read in the same order and adds one step: the model
whose plans cross its finish line most often is the one the projector and then the controller comparison
of \autoref{sec:res:uav:controller} are applied to, the one place in this thesis where MuJoCo MPC replaces
the tracking controller (\autoref{sec:res:uav:scurve}). Each scene closes on its own conclusion.

UAV-corridor and UAV-s-curve use one training seed, and every generative model there is a temporal U-Net
with $3.96$--$3.97$\,M parameters, trained with the same optimiser settings (\autoref{tab:train});
UAV-pillars flies the models of D3IL-avoiding at its five training seeds (\autoref{sec:setup:protocol:uav}).
Every flight is tracked by the cascaded geometric controller of \autoref{sec:method:geometric} unless
stated otherwise. The diffusion
model is trained at $\nfe=20$ and is compared at the budget it was trained with.

\subsection{UAV-pillars}\label{sec:res:uav:pillars}

This scene is D3IL-avoiding flown. The arena is the obstacle field scaled by $36$, the constraint sets are
the three geometries of the table mapped into it, and the planner is the one of \autoref{sec:res:state},
not retrained (\autoref{sec:setup:tasks:uav}). What is read here is therefore not a new benchmark but a
retest of the result of D3IL-avoiding on another plant. The configurations flown are the Pareto-optimal ones
of that benchmark, the two average-velocity models at one and two network evaluations on its cost frontier
(\autoref{sec:res:avoiding:pareto}, \autoref{tab:avoiding-conclusion}), and they are flown again, before and
after per-step projection, to test what else happens when a quadrotor and its tracking controller execute the
same planner configurations instead of the manipulator: whether the result holds, and where the flown outcome departs from the
table's. Each flown outcome is set beside its table outcome at the same protocol, five training seeds $\times$
two episodes per geometry, averaged over the three geometries. Every flight of the scene ended by crossing the goal line
or by touching a pillar; none reached the episode limit and none was aborted.

\subsubsection{From D3IL-avoiding to UAV-pillars}\label{sec:res:uav:models}

Before projection the vehicle changes little (\autoref{tab:uav-pillars-raw}, upper block). The four
unprojected cells read $0.033$ success with constraint satisfaction on D3IL-avoiding and $0.033$ to
$0.100$ on UAV-pillars, with the same $15$ to $20$ violating control steps per episode and the same sixty
control steps to the goal line; the planner's time per action is the same to within a millisecond, since
nothing about the plant enters it. Nine of the $120$ unprojected flights end in a collision with a pillar,
and the rest cross the goal line.

\begin{table}[htpb]
  \caption[UAV-pillars: the avoiding configurations on D3IL-avoiding and on UAV-pillars]{The two
  average-velocity models (CI-MeanFM at $\alpha_{\mathrm{end}}=0.2$) at $\nfe=1$ and $2$, before projection
  and after per-step projection under temporal consistency, on tightened constraints: executed by the
  manipulator on D3IL-avoiding and flown by the quadrotor on UAV-pillars, one row each. Columns as in
  \autoref{tab:avoiding-raw-models}. The protocol of \ac{DPCC} on both: five training seeds $\times$ two
  episodes per geometry, averaged over the three geometries; mean $\pm$ the spread over the seeds. Every
  UAV-pillars flight without success ended in a collision with a pillar. \selectedmark\ the selected
  configuration of D3IL-avoiding.}
  \label{tab:uav-pillars-raw}
  \centering
  \footnotesize
  \setlength{\tabcolsep}{4pt}
  \begin{tabularx}{\linewidth}{@{}L r l rrrr@{}}
    \toprule
    Model & $\nfe$ & Environment & Success & S\&C & \shortstack{Violating\\steps} & Steps \\
    \midrule
    \multicolumn{7}{@{}l}{\emph{before projection}}\\
    MeanFM & 1 & D3IL-avoiding & $1.000 \pm 0.000$ & $0.033 \pm 0.075$ & $15.5 \pm 2.1$ & $63.4 \pm 5.6$ \\
              &   & UAV-pillars & $1.000 \pm 0.000$ & $0.100 \pm 0.091$ & $19.9 \pm 0.6$ & $62.7 \pm 3.3$ \\
    \addlinespace[2pt]
    MeanFM & 2 & D3IL-avoiding & $1.000 \pm 0.000$ & $0.033 \pm 0.075$ & $15.5 \pm 0.4$ & $62.7 \pm 1.8$ \\
              &   & UAV-pillars & $0.800 \pm 0.274$ & $0.100 \pm 0.091$ & $19.1 \pm 1.8$ & $60.9 \pm 3.2$ \\
    \addlinespace[2pt]
    CI-MeanFM & 1 & D3IL-avoiding & $1.000 \pm 0.000$ & $0.033 \pm 0.075$ & $15.6 \pm 0.7$ & $60.5 \pm 1.5$ \\
              &   & UAV-pillars & $1.000 \pm 0.000$ & $0.100 \pm 0.149$ & $18.1 \pm 3.2$ & $59.7 \pm 0.7$ \\
    \addlinespace[2pt]
    CI-MeanFM & 2 & D3IL-avoiding & $1.000 \pm 0.000$ & $0.033 \pm 0.075$ & $15.7 \pm 0.7$ & $62.9 \pm 1.2$ \\
              &   & UAV-pillars & $0.900 \pm 0.224$ & $0.033 \pm 0.075$ & $19.8 \pm 1.7$ & $62.3 \pm 4.2$ \\
    \midrule
    \multicolumn{7}{@{}l}{\emph{after per-step projection, temporal consistency}}\\
    \mainconfig{\selectedmark\ MeanFM} & \mainconfig{1} & \mainconfig{D3IL-avoiding} & \mainconfig{$0.967 \pm 0.075$} & \mainconfig{$0.967 \pm 0.075$} & \mainconfig{$0.0$} & \mainconfig{$58.6 \pm 2.2$} \\
              &   & \mainconfig{UAV-pillars} & \mainconfig{$0.600 \pm 0.190$} & \mainconfig{$0.600 \pm 0.190$} & \mainconfig{$0.0$} & \mainconfig{$61.4 \pm 2.3$} \\
    \addlinespace[2pt]
    MeanFM & 2 & D3IL-avoiding & $0.967 \pm 0.075$ & $0.967 \pm 0.075$ & $0.0$ & $59.4 \pm 2.0$ \\
              &   & UAV-pillars & $0.667 \pm 0.236$ & $0.667 \pm 0.236$ & $0.0$ & $60.2 \pm 1.6$ \\
    \addlinespace[2pt]
    CI-MeanFM & 1 & D3IL-avoiding & $1.000 \pm 0.000$ & $1.000 \pm 0.000$ & $0.0$ & $59.2 \pm 1.2$ \\
              &   & UAV-pillars & $0.633 \pm 0.217$ & $0.633 \pm 0.217$ & $0.0$ & $60.7 \pm 1.6$ \\
    \addlinespace[2pt]
    CI-MeanFM & 2 & D3IL-avoiding & $1.000 \pm 0.000$ & $1.000 \pm 0.000$ & $0.0$ & $60.1 \pm 3.4$ \\
              &   & UAV-pillars & $0.700 \pm 0.139$ & $0.700 \pm 0.139$ & $0.0$ & $60.6 \pm 0.6$ \\
    \bottomrule
  \end{tabularx}
\end{table}

After projection the declared constraints hold on UAV-pillars as they hold on D3IL-avoiding, and a third
of the flights still fail. The projected rows read $0.0$ violating control steps on every flight that
reached the goal line --- $78$ flights, none of which crosses a halfspace or enters the keep-out disk ---
so what the per-step projector delivers on D3IL-avoiding it delivered on every one of the $120$ projected
flights, up to the moment each ended. Yet success falls
from $0.967$ and $1.000$ to $0.600$ to $0.700$: $42$ of the $120$ projected flights end in a collision with
a pillar, with no violation of the declared constraint set up to that moment. Since no projected flight
violates a constraint, success with constraint satisfaction falls with success. The loss sits on one
geometry. On \emph{both-hard} two to four of ten projected flights succeed and the other six to eight end
in a collision, while on the two single-halfspace geometries seven to nine of ten succeed
(\autoref{tab:uav-pillars-geo}). The two models lose the same amount on the same geometry at both budgets,
so nothing here re-orders them.

\begin{table}[htpb]
  \caption[UAV-pillars: the projected flights per geometry]{The projected rows of
  \autoref{tab:uav-pillars-raw} per constraint geometry: success with constraint satisfaction on
  D3IL-avoiding and on UAV-pillars. Every UAV-pillars flight counted as a failure here ended in a
  collision with a pillar.}
  \label{tab:uav-pillars-geo}
  \centering
  \small
  \begin{tabularx}{\linewidth}{@{}L r rr rr rr@{}}
    \toprule
    & & \multicolumn{2}{c}{\emph{top-right-hard}} & \multicolumn{2}{c}{\emph{top-left-hard}} & \multicolumn{2}{c}{\emph{both-hard}} \\
    \cmidrule(lr){3-4}\cmidrule(lr){5-6}\cmidrule(l){7-8}
    Model & $\nfe$ & \shortstack{D3IL-\\avoiding} & \shortstack{UAV-\\pillars} & \shortstack{D3IL-\\avoiding} & \shortstack{UAV-\\pillars} & \shortstack{D3IL-\\avoiding} & \shortstack{UAV-\\pillars} \\
    \midrule
    MeanFM    & 1 & 0.90 & 0.90 & 1.00 & 0.70 & 1.00 & 0.20 \\
    MeanFM    & 2 & 0.90 & 0.80 & 1.00 & 0.80 & 1.00 & 0.40 \\
    CI-MeanFM & 1 & 1.00 & 0.80 & 1.00 & 0.90 & 1.00 & 0.20 \\
    CI-MeanFM & 2 & 1.00 & 0.90 & 1.00 & 0.90 & 1.00 & 0.30 \\
    \bottomrule
  \end{tabularx}
\end{table}

Two things about these collisions are the quadrotor's own: why they happen, and what follows from them.
They happen because of the plans and the scene, not the tracking. Five of the six pillars are in no
constraint set a geometry declares --- one or two halfspaces and one keep-out disk, which covers the sixth
(\autoref{sec:setup:tasks:avoiding}) --- so the projector keeps the plan away from none of them, and the
demonstrations of D3IL-avoiding pass them at the rod's own radius, which that benchmark never scores. The
map of the scene sends that radius onto the rotor reach (\autoref{sec:setup:tasks:uav}), so a commanded path
that skims a post on D3IL-avoiding brings the rotors into the pillar on UAV-pillars: on $50$ of the $51$
flights that ended in a collision the commanded path itself passes within the reach of that pillar, and only
one collision is decided by the tracking lag. On \emph{both-hard} the two halfspaces funnel the plan through
the centre of the field and past the outer pillars of the third row, which is where $29$ of the $42$
projected collisions are; on the single-halfspace geometries the collisions are on the second row, where the
plan rounds the pillar next to the keep-out disk.

What follows from a collision is where the quadrotor differs from the manipulator in kind. The arm stands on
a fixed base and holds its end effector under position control, so an end effector that meets an obstacle is
stopped and stays under control. The quadrotor is held in the air by its rotors alone: a rotor that strikes a
pillar no longer gives the thrust the attitude loop relies on, and in most cases the vehicle loses control of
its attitude and falls. A collision is therefore the end of a flight, not an event to be counted within it.
The simulation does not model the strike itself; the scene ends the flight at its first contact with a
pillar, and the tables read it through the success column (\autoref{sec:setup:metrics:uav}).

\FloatBarrier
\subsubsection{Executed Paths: the Manipulator and the Quadrotor}\label{sec:res:uav:pillars:projected}

\autoref{fig:uav-pillars-paths} draws the flights of the selected configuration, MeanFM at $\nfe=1$, on
each geometry, before and after projection. Before projection every flight crosses the goal line and
almost every one crosses the declared boundary on its way: on \emph{top-right-hard} the flights that
take the upper passage cut through the halfspace and those that take the lower one pass through the
keep-out disk, and on \emph{both-hard} both passages cross a halfspace. After projection no path enters
the shaded region or the disk; the paths hug the dashed tightened boundary instead, and where a hugging
path meets a pillar that stands at the boundary the flight ends in a collision --- one flight on
\emph{top-right-hard}, three on \emph{top-left-hard}, eight on \emph{both-hard}, at the third-row pillars
on either side of the finish line's centre. The paths that do reach the line are the paths of
D3IL-avoiding (\autoref{fig:avoiding-paths}) at scale $36$.

\begin{figure}[htpb]
  \centering
  \includegraphics[width=\linewidth]{fig_uav_pillars_paths}
  \caption[UAV-pillars: the flown paths of the selected configuration]{UAV-pillars from above, MeanFM at
  $\nfe=1$: the ten flights of each constraint geometry (columns) before projection (upper row) and under
  per-step projection with temporal consistency (lower row), one seed per pair of flights. Shaded
  regions and the dashed tightened boundary as in \autoref{fig:constraints-uav}; the six pillars in red.
  Green: violation-free; red: violated a constraint; dotted: no success; a cross marks a collision with a
  pillar, which ends the flight.}
  \label{fig:uav-pillars-paths}
\end{figure}

\FloatBarrier
\subsubsection{Conclusion}\label{sec:res:uav:pillars:conclusion}

UAV-pillars flies the Pareto-optimal operating points of D3IL-avoiding again, on the quadrotor, to see what
else happens when the plant changes, and it reaffirms D3IL-avoiding. Its demonstrations, planner, projection
and constraint rows are those of D3IL-avoiding, and only the plant that executes each commanded position is
replaced (\autoref{sec:method:env:uav}). As far as the declared constraints go, the result of D3IL-avoiding holds when
flown: before projection the flown plans violate them as the executed ones do, after projection neither
violates them at all, at the same control steps to within three and the same time per action to within a
millisecond. The two average-velocity models keep their order and their gap, CI-MeanFM succeeding within the
constraints on one flight in thirty more than analytic average-velocity matching at both budgets, on
UAV-pillars as on D3IL-avoiding, and per-step projection is again what turns plans that violate the declared
constraints in almost every episode into plans that violate none.

Where the flown outcome departs from the table, the cause is the plant and its controller, and what they
change against the manipulator is twofold. The vehicle tracks more
closely: the $95$th percentile of a flight's gap to the commanded position, averaged over the flights, is
$0.007$ in the table's units, against the
manipulator's own $0.02$ to $0.04$ and well inside the tightening margin of $0.025$, so the flown path keeps
the constraint satisfaction the projection gives the plan. And an obstacle costs more. Success within the
constraints falls to $0.600$--$0.700$ against $0.967$--$1.000$, entirely through collisions with the five
pillars the constraint sets leave out: the plans pass them as closely as the demonstrations passed the posts,
at the rod's radius, which the scale of the scene makes the rotor reach --- on $50$ of the $51$ collisions the
commanded path itself is inside the reach --- and for a quadrotor a collision is catastrophic, a rotor strike
that in most cases ends in lost control, so the flight ends at the first contact. Both models lose the same
share on the same geometry, so the collisions do not re-order them. The declared constraint set transfers
exactly; freedom from collision with the obstacles it leaves out does not.
Only the Pareto-optimal operating points of D3IL-avoiding were flown on this scene, since it tests what
the change of plant does to them; the diffusion baseline and instantaneous-velocity matching were not flown,
and the dominance over the baseline remains a result of D3IL-avoiding.

\subsection{UAV-corridor}\label{sec:res:uav:corridor}

The two constraints are read one at a time, as the avoiding geometries are: every table of this scene has a
\emph{tilt} block and a \emph{hump} block, and nothing is averaged across them. A flight succeeds when it
reaches the end of the corridor, the scene's finish line, and since almost no flight of this scene is free
of violations, the constraint is read as the number of violating control steps per flight
(\autoref{sec:setup:metrics:uav}). Every value is the mean over the ten flights of a cell with the sample
standard deviation over them.

\subsubsection{Generative Models}\label{sec:res:uav:corridor:raw}
Without projection every model reaches the end of the corridor on every flight at every budget, and every
flight violates the constraint (\autoref{tab:uav-corridor-raw}): a flight violates the tilt on $74.5$ to
$99.4$ control steps on average and the hump on $28.6$ to $36.6$. The plans follow the level lanes of the
demonstrations, which were flown before either boundary existed (\autoref{sec:setup:data:uav}), so the two
blocks are the same flights read against two constraints: they share their steps, their time per action
differs by at most $1.5$\,ms between the two runs, and they differ in what they violate. The flow-based
models lie within one standard deviation of one another in violating steps; the diffusion baseline has the
most under the tilt and the fewest over the hump. The flights take $242$ to $251$ control steps to the end
of the corridor, and one action costs $9.1$ to $9.4$\,ms at one evaluation, $170.8$ to $172.3$\,ms at
twenty and about $178$\,ms for the baseline. As before projection on D3IL-avoiding, every configuration
succeeds, and no cost frontier is drawn for it.

\begin{table}[htpb]
  \caption[UAV-corridor: the generative models before projection]{UAV-corridor with the test-time
  constraints, \textbf{before projection}: the plan flown as generated, ten flights per row spread over the
  three routes, one training seed, one block per constraint. The grid: the two average-velocity models at
  $\nfe\in\{1,2,3\}$, instantaneous-velocity matching at $\nfe\in\{1,2,3,5,20\}$, the diffusion baseline at
  its trained budget. Each value mean $\pm$ the sample standard deviation over the flights; \emph{Success}:
  reaches the end of the corridor (\autoref{sec:setup:metrics:uav}); \emph{Violating steps}: control steps
  of the flight that violate the constraint; \emph{Steps}: control steps to the end of the corridor.
  \baselinemark\ the diffusion model of \ac{DPCC}.}
  \label{tab:uav-corridor-raw}
  \centering
  \footnotesize
  \setlength{\tabcolsep}{4pt}
  \begin{tabularx}{\linewidth}{@{}L r rrrr@{}}
    \toprule
    Model & $\nfe$ & Success & \shortstack{Violating\\steps} & Steps & ms/step \\
    \midrule
    \multicolumn{6}{@{}l}{\emph{tilt}}\\
    MeanFM & 1 & $1.000 \pm 0.000$ & $91.3 \pm 34.8$ & $246.4 \pm 7.0$ & $9.4 \pm 0.4$ \\
    MeanFM & 2 & $1.000 \pm 0.000$ & $92.9 \pm 30.2$ & $247.3 \pm 4.8$ & $18.1 \pm 0.4$ \\
    MeanFM & 3 & $1.000 \pm 0.000$ & $92.0 \pm 30.4$ & $246.5 \pm 5.0$ & $27.0 \pm 0.4$ \\
    CI-MeanFM, $\alpha_{\mathrm{end}}=0.2$ & 1 & $1.000 \pm 0.000$ & $90.0 \pm 32.6$ & $242.4 \pm 4.3$ & $9.3 \pm 0.4$ \\
    CI-MeanFM, $\alpha_{\mathrm{end}}=0.2$ & 2 & $1.000 \pm 0.000$ & $93.8 \pm 30.0$ & $246.2 \pm 5.5$ & $18.4 \pm 0.4$ \\
    CI-MeanFM, $\alpha_{\mathrm{end}}=0.2$ & 3 & $1.000 \pm 0.000$ & $93.1 \pm 30.1$ & $245.4 \pm 4.9$ & $27.3 \pm 0.4$ \\
    FM & 1 & $1.000 \pm 0.000$ & $74.5 \pm 25.6$ & $250.9 \pm 5.2$ & $9.1 \pm 0.4$ \\
    FM & 2 & $1.000 \pm 0.000$ & $83.9 \pm 28.4$ & $251.2 \pm 2.9$ & $17.7 \pm 0.4$ \\
    FM & 3 & $1.000 \pm 0.000$ & $87.0 \pm 28.3$ & $251.4 \pm 2.7$ & $25.9 \pm 0.5$ \\
    FM & 5 & $1.000 \pm 0.000$ & $89.9 \pm 28.4$ & $249.7 \pm 2.5$ & $43.4 \pm 0.4$ \\
    FM & 20 & $1.000 \pm 0.000$ & $94.7 \pm 27.0$ & $249.5 \pm 2.9$ & $170.8 \pm 0.4$ \\
    \mainconfig{\baselinemark\ Diffusion} & \mainconfig{20} & \mainconfig{$1.000 \pm 0.000$} & \mainconfig{$99.4 \pm 27.0$} & \mainconfig{$242.4 \pm 4.2$} & \mainconfig{$178.1 \pm 0.8$} \\
    \midrule
    \multicolumn{6}{@{}l}{\emph{hump}}\\
    MeanFM & 1 & $1.000 \pm 0.000$ & $36.0 \pm 8.9$ & $246.4 \pm 7.0$ & $9.4 \pm 0.4$ \\
    MeanFM & 2 & $1.000 \pm 0.000$ & $35.6 \pm 8.1$ & $247.3 \pm 4.8$ & $18.2 \pm 0.4$ \\
    MeanFM & 3 & $1.000 \pm 0.000$ & $35.7 \pm 7.8$ & $246.5 \pm 5.0$ & $27.2 \pm 0.4$ \\
    CI-MeanFM, $\alpha_{\mathrm{end}}=0.2$ & 1 & $1.000 \pm 0.000$ & $36.6 \pm 8.4$ & $242.4 \pm 4.3$ & $9.3 \pm 0.4$ \\
    CI-MeanFM, $\alpha_{\mathrm{end}}=0.2$ & 2 & $1.000 \pm 0.000$ & $35.6 \pm 7.9$ & $246.2 \pm 5.5$ & $18.0 \pm 0.4$ \\
    CI-MeanFM, $\alpha_{\mathrm{end}}=0.2$ & 3 & $1.000 \pm 0.000$ & $35.6 \pm 7.6$ & $245.4 \pm 4.9$ & $27.1 \pm 0.4$ \\
    FM & 1 & $1.000 \pm 0.000$ & $34.7 \pm 6.8$ & $250.9 \pm 5.2$ & $9.1 \pm 0.4$ \\
    FM & 2 & $1.000 \pm 0.000$ & $32.6 \pm 7.8$ & $251.2 \pm 2.9$ & $17.5 \pm 0.3$ \\
    FM & 3 & $1.000 \pm 0.000$ & $32.6 \pm 8.0$ & $251.4 \pm 2.7$ & $26.3 \pm 0.4$ \\
    FM & 5 & $1.000 \pm 0.000$ & $32.0 \pm 7.9$ & $249.7 \pm 2.5$ & $43.6 \pm 0.4$ \\
    FM & 20 & $1.000 \pm 0.000$ & $31.3 \pm 7.6$ & $249.5 \pm 2.9$ & $172.3 \pm 2.6$ \\
    \mainconfig{\baselinemark\ Diffusion} & \mainconfig{20} & \mainconfig{$1.000 \pm 0.000$} & \mainconfig{$28.6 \pm 7.7$} & \mainconfig{$242.4 \pm 4.2$} & \mainconfig{$178.2 \pm 2.7$} \\
    \bottomrule
  \end{tabularx}
\end{table}

\FloatBarrier
\subsubsection{Projection}\label{sec:res:uav:projected}
After projection every flow-based configuration still reaches the end of the corridor on every flight,
except per-step projection at one and two evaluations over the hump, whose flights stop short of the apex
(\autoref{tab:uav-corridor}). Projection removes most of the violating steps and rarely all of them: a
flight's count falls from $74.5$ to $99.4$ under the tilt and from $28.6$ to $36.6$ over the hump to $1.4$
to $10.4$ and $0.2$ to $8.3$ for the flow-based models, and to $0.8$ and $0.0$ for the diffusion baseline
at its best rule, at $666$ and $654$\,ms per action against $26$ to $417$\,ms for the flow-based models.
At the route's goal point, the line these flights were first scored at, only the baseline's success
differs, $0.000$ under the tilt and $0.300$ over the hump at its best rule: $41$ of its $60$ projected
flights had passed the end of the corridor and were still flying towards the goal when the episode limit
ended them (\autoref{sec:setup:metrics:uav}). The budgets are read as one ladder across the scene: the
per-step projector is active at every budget, endpoint projection at $\nfe\ge 3$, where it has a guiding
step (\autoref{sec:res:constraints:degenerate}), so at $\nfe=1$ and $\nfe=2$ the table carries per-step
rows only.

\begin{table}[htpb]
  \caption[UAV-corridor: the generative models after projection]{UAV-corridor \textbf{after projection}:
  each model at each budget under the per-step projector of \ac{DPCC} and, from $\nfe=3$, under endpoint
  projection, at its best selection rule (the order and every rule in \autoref{app:uav-corridor-rules});
  tightened constraints, activation threshold $0.5$, ten flights per cell, one training seed. Columns as in
  \autoref{tab:uav-corridor-raw}; \emph{---}: no flight reaches the end of the corridor.
  \baselinemark\ the diffusion model of \ac{DPCC}.}
  \label{tab:uav-corridor}
  \centering
  \footnotesize
  \setlength{\tabcolsep}{4pt}
  \begin{tabularx}{\linewidth}{@{}L r c rrrr@{}}
    \toprule
    Model & $\nfe$ & Rule & Success & \shortstack{Violating\\steps} & Steps & ms/step \\
    \midrule
    \multicolumn{7}{@{}l}{\emph{tilt, per-step projection}}\\
    MeanFM & 1 & $c$ & $1.000 \pm 0.000$ & $10.2 \pm 1.1$ & $253.2 \pm 4.4$ & $27.4 \pm 0.5$ \\
    MeanFM & 2 & $c$ & $1.000 \pm 0.000$ & $10.3 \pm 0.9$ & $253.9 \pm 4.2$ & $36.4 \pm 0.8$ \\
    MeanFM & 3 & $c$ & $1.000 \pm 0.000$ & $6.6 \pm 0.8$ & $256.4 \pm 2.9$ & $93.6 \pm 1.1$ \\
    CI-MeanFM, $\alpha_{\mathrm{end}}=0.2$ & 1 & $c$ & $1.000 \pm 0.000$ & $10.3 \pm 0.9$ & $252.0 \pm 3.7$ & $27.9 \pm 0.5$ \\
    CI-MeanFM, $\alpha_{\mathrm{end}}=0.2$ & 2 & $c$ & $1.000 \pm 0.000$ & $10.4 \pm 0.7$ & $251.9 \pm 2.4$ & $36.1 \pm 0.3$ \\
    CI-MeanFM, $\alpha_{\mathrm{end}}=0.2$ & 3 & $c$ & $1.000 \pm 0.000$ & $5.0 \pm 0.8$ & $254.0 \pm 1.8$ & $93.4 \pm 1.2$ \\
    FM & 1 & $c$ & $1.000 \pm 0.000$ & $10.2 \pm 1.5$ & $259.8 \pm 9.2$ & $27.6 \pm 1.5$ \\
    FM & 2 & $c$ & $1.000 \pm 0.000$ & $9.8 \pm 1.4$ & $253.9 \pm 9.4$ & $34.7 \pm 0.5$ \\
    FM & 3 & $c$ & $1.000 \pm 0.000$ & $5.8 \pm 1.5$ & $247.2 \pm 3.4$ & $79.2 \pm 3.2$ \\
    FM & 5 & $c$ & $1.000 \pm 0.000$ & $4.4 \pm 1.5$ & $247.6 \pm 2.0$ & $116.9 \pm 2.4$ \\
    FM & 20 & $c$ & $1.000 \pm 0.000$ & $2.1 \pm 0.7$ & $252.6 \pm 2.7$ & $360.4 \pm 1.5$ \\
    \mainconfig{\baselinemark\ Diffusion} & \mainconfig{20} & \mainconfig{$c$} & \mainconfig{$1.000 \pm 0.000$} & \mainconfig{$0.8 \pm 1.8$} & \mainconfig{$290.3 \pm 15.5$} & \mainconfig{$666.3 \pm 16.7$} \\
    \addlinespace[3pt]
    \multicolumn{7}{@{}l}{\emph{tilt, endpoint projection}}\\
    MeanFM & 3 & $c$ & $1.000 \pm 0.000$ & $7.6 \pm 1.3$ & $255.2 \pm 4.5$ & $73.5 \pm 0.8$ \\
    CI-MeanFM, $\alpha_{\mathrm{end}}=0.2$ & 3 & $c$ & $1.000 \pm 0.000$ & $7.1 \pm 1.0$ & $253.2 \pm 3.5$ & $72.7 \pm 0.5$ \\
    FM & 3 & $c$ & $1.000 \pm 0.000$ & $6.9 \pm 1.8$ & $257.8 \pm 3.3$ & $71.6 \pm 0.6$ \\
    FM & 5 & $c$ & $1.000 \pm 0.000$ & $4.9 \pm 1.9$ & $256.6 \pm 2.9$ & $115.2 \pm 1.4$ \\
    FM & 20 & $c$ & $1.000 \pm 0.000$ & $1.4 \pm 0.7$ & $269.8 \pm 10.5$ & $416.9 \pm 0.8$ \\
    \midrule
    \multicolumn{7}{@{}l}{\emph{hump, per-step projection}}\\
    MeanFM & 1 & $c$ & $0.000 \pm 0.000$ & $0.0 \pm 0.0$ & --- & $26.4 \pm 0.0$ \\
    MeanFM & 2 & $t$ & $0.100 \pm 0.316$ & $17.9 \pm 30.0$ & $387.0$ & $42.3 \pm 10.9$ \\
    MeanFM & 3 & $c$ & $1.000 \pm 0.000$ & $4.9 \pm 0.9$ & $286.3 \pm 4.3$ & $91.1 \pm 1.8$ \\
    CI-MeanFM, $\alpha_{\mathrm{end}}=0.2$ & 1 & $t$ & $0.000 \pm 0.000$ & $0.0 \pm 0.0$ & --- & $26.6 \pm 0.3$ \\
    CI-MeanFM, $\alpha_{\mathrm{end}}=0.2$ & 2 & $t$ & $0.200 \pm 0.422$ & $19.8 \pm 32.1$ & $388.0 \pm 2.8$ & $43.0 \pm 12.2$ \\
    CI-MeanFM, $\alpha_{\mathrm{end}}=0.2$ & 3 & $c$ & $1.000 \pm 0.000$ & $4.6 \pm 0.8$ & $282.7 \pm 3.6$ & $89.2 \pm 1.9$ \\
    FM & 1 & $c$ & $0.000 \pm 0.000$ & $0.0 \pm 0.0$ & --- & $26.2 \pm 0.2$ \\
    FM & 2 & $r$ & $0.000 \pm 0.000$ & $0.0 \pm 0.0$ & --- & $34.5 \pm 0.4$ \\
    FM & 3 & $c$ & $1.000 \pm 0.000$ & $5.6 \pm 0.5$ & $277.0 \pm 7.3$ & $77.1 \pm 2.7$ \\
    FM & 5 & $c$ & $1.000 \pm 0.000$ & $3.9 \pm 0.7$ & $277.8 \pm 7.0$ & $114.1 \pm 2.3$ \\
    FM & 20 & $c$ & $1.000 \pm 0.000$ & $1.8 \pm 0.4$ & $285.1 \pm 5.6$ & $355.2 \pm 2.2$ \\
    \mainconfig{\baselinemark\ Diffusion} & \mainconfig{20} & \mainconfig{$r$} & \mainconfig{$1.000 \pm 0.000$} & \mainconfig{$0.0 \pm 0.0$} & \mainconfig{$320.8 \pm 53.6$} & \mainconfig{$654.2 \pm 89.7$} \\
    \addlinespace[3pt]
    \multicolumn{7}{@{}l}{\emph{hump, endpoint projection}}\\
    MeanFM & 3 & $c$ & $1.000 \pm 0.000$ & $8.0 \pm 1.5$ & $287.7 \pm 4.9$ & $72.1 \pm 0.3$ \\
    CI-MeanFM, $\alpha_{\mathrm{end}}=0.2$ & 3 & $c$ & $1.000 \pm 0.000$ & $8.0 \pm 1.2$ & $284.6 \pm 2.8$ & $71.9 \pm 0.6$ \\
    FM & 3 & $c$ & $1.000 \pm 0.000$ & $8.3 \pm 1.1$ & $299.3 \pm 4.7$ & $71.9 \pm 1.0$ \\
    FM & 5 & $c$ & $1.000 \pm 0.000$ & $3.9 \pm 0.3$ & $292.8 \pm 3.7$ & $112.4 \pm 0.6$ \\
    FM & 20 & single & $1.000 \pm 0.000$ & $0.2 \pm 0.4$ & $316.2 \pm 6.3$ & $291.0 \pm 0.8$ \\
    \bottomrule
  \end{tabularx}
\end{table}

The budget makes the violations fewer and shallower and does not remove them. Wherever the flights reach
the end of the corridor, the violating steps per flight fall with the budget, under both constraints and
both projection methods. Under the tilt they fall from $9.8$ to $11.1$ at one and two evaluations to $2.1$
to $3.2$ at twenty under per-step projection and to $1.4$ to $2.5$ under endpoint projection, and the
deepest of them from $8.8$ to $11.9$\,cm to $1.3$ to $2.2$ and $1.0$ to $1.3$\,cm. Over the hump they fall
from $4.6$ to $6.5$ under per-step and $8.0$ to $10.4$ under endpoint projection at three evaluations to
$1.8$ to $2.2$ and $0.2$ to $0.7$ at twenty. The budget at which a flight stops reaching the finish line lies
at the other end of the ladder: under per-step projection at one and two evaluations $177$ of the $180$
flow-based flights over the hump stop short of the apex, and the other three reach the end of the corridor
only after drifting across the boundary of its wall. The baseline, at its one budget, is the slowest
planner on the scene: its projected flights need $290$ (tilt) and $321$ (hump) control steps to the end of
the corridor, against $247$ to $270$ and $277$ to $316$ for the flow-based ones, and $623$ to $680$\,ms
per action, the highest price of the scene.

The cost frontier of this scene is drawn after projection only, on violating steps per flight against the
time to compute one action (\autoref{fig:uav-corridor-tradeoff}, top). A configuration takes part when every one
of its flights reaches the end of the corridor --- its success within $0.05$ of the best, as on
D3IL-avoiding --- and is on the frontier when no other configuration leaves fewer violating steps at no more
time. Both frontiers are staircases along the budget, and the budget is the trade between the two axes:
every step up in network evaluations buys fewer violating steps with more time per action. Under the tilt a
flight violates on about ten control steps at one and two evaluations ($27$ to $36$\,ms per action), on four
to eight at three and five ($72$ to $117$\,ms) and on one to two at twenty ($360$ to $417$\,ms); over the hump
on about eight at three evaluations under endpoint projection and four to six under per-step projection
($72$ to $91$\,ms), about four at five ($112$ to $114$\,ms) and under two at twenty ($291$ to $355$\,ms). Under the tilt it runs from analytic average-velocity matching at one evaluation, $10.2$ violating
steps at $27.4$\,ms, through instantaneous-velocity matching at two, three, five and twenty evaluations and
consistency-interpolated average-velocity matching at three, to the baseline, $0.8$ at $666.3$\,ms. Over the
hump it runs from consistency-interpolated average-velocity matching at three evaluations under endpoint
projection, $8.0$ at $71.9$\,ms, through instantaneous-velocity matching and the consistency-interpolated
model at three evaluations under per-step projection and instantaneous-velocity matching at five under
endpoint projection, to instantaneous-velocity matching at twenty with a single candidate under endpoint
projection, $0.2$ at $291.0$\,ms, and the baseline, $0.0$ at $654.2$\,ms.
No model is ahead at every budget. Where the three flow-based models share a budget, the average-velocity
objectives are on the frontier at its cheap end: under the tilt at one evaluation analytic average-velocity
matching leaves the same $10.2$ violating steps as instantaneous-velocity matching $0.2$\,ms sooner, and at
three evaluations the consistency-interpolated model leaves the fewest under per-step projection on both
constraints, $5.0$ and $4.6$ against $5.8$ and $5.6$, at about a sixth more time, and over the hump under
endpoint projection $8.0$ against $8.3$ at the same $71.9$\,ms. Each of these differences is about one
standard deviation of the flights. Per-step projection at one and two evaluations over the hump is off the
frontier: its flights stop before the apex, most of them with no violating step.

The configurations on that frontier fall into groups that violate about equally, and within a group the
comparison is the one of D3IL-avoiding: which of them reaches the finish line in fewer control steps at no
more time (\autoref{fig:uav-corridor-tradeoff}, bottom). The groups are where the violating steps jump by
more than two: under the tilt about ten (one and two evaluations), four to eight (three and five) and up
to two (twenty, and the baseline); over the hump about eight (endpoint projection at three evaluations), four
to six (per-step projection at three, and five) and up to two (twenty, and the baseline). Among the one- and
two-evaluation configurations under the tilt the average-velocity objectives reach the line first: analytic
average-velocity matching at one evaluation in $253.2$ control steps at $27.4$\,ms, against $259.8$ at
$27.6$\,ms for instantaneous-velocity matching, and the consistency-interpolated model in $252.0$ and $251.9$
steps at one and two evaluations. In the middle group under the tilt instantaneous-velocity matching under
per-step projection at three evaluations is the shortest, $247.2$ steps at $79.2$\,ms, and the endpoint
cells of it and of the consistency-interpolated model are about $7$\,ms cheaper at $253.2$ to $257.8$ steps.
Over the hump the consistency-interpolated model under endpoint projection dominates its group, $284.6$ steps
against $287.7$ and $299.3$ at $72$\,ms, and instantaneous-velocity matching under per-step projection at
three evaluations dominates the next, $277.0$ steps at $77.1$\,ms. Among the configurations with the fewest
violating steps, the band of at most two, instantaneous-velocity matching at twenty evaluations is on the
frontier and the baseline is not, on steps and time within that band --- in which the baseline itself
violates least, $0.8$ and $0.0$: under the tilt per-step projection, $252.6$ steps at $360$\,ms, dominates
endpoint projection ($269.8$, $417$\,ms) and the baseline ($290.3$, $666$\,ms) on those two axes; over the
hump endpoint projection with one candidate is the cheaper, $316.2$ steps at $291$\,ms, per-step projection
the shorter, $285.1$ at $355$\,ms, and the baseline, $320.8$ steps at $654$\,ms, is dominated on both. Most of these differences are about one standard deviation of the
flights; the largest, over the hump between the consistency-interpolated model and instantaneous-velocity
matching under endpoint projection, is three.

\begin{figure}[htpb]
  \centering
  \includegraphics[width=\linewidth]{fig_uav_corridor_tradeoff}
  \caption[UAV-corridor after projection: violating steps and control steps against cost]{UAV-corridor after
  projection, one point per model, budget and projection method at its best rule (\autoref{tab:uav-corridor});
  colour the model, circles per-step projection, squares endpoint projection, the number its budget.
  (a, b) Violating steps per flight against the time to compute one action; rings and the dashed staircase
  mark the non-dominated points among those reaching the end of the corridor on every flight (hollow: not).
  Bands I--III group configurations separated by more than two violating steps. (c, d) Within each group,
  control steps to the finish line against time, its non-dominated points ringed. Ten flights per point,
  one training seed.}
  \label{fig:uav-corridor-tradeoff}
\end{figure}

Where the violating steps lie is seen in the flights themselves (\autoref{fig:uav-corridor-side},
\autoref{app:uav-corridor-side}), drawn
for instantaneous-velocity matching, the one flow-based model flown at twenty evaluations, without
projection and under per-step projection at one and at twenty evaluations. Under the tilt the unprojected
flights stay level, so their margin to the leaned plane shrinks steadily along the corridor, to as much
as $0.37$\,m inside the plane at its end. The projected flights move sideways and down until they hold a few
centimetres clear of the plane, and every one of them cuts into it in the last $0.3$\,m of the plane's
span: every violating step of every projected flight, on every model, budget and projection method, lies
between $x=1.71$ and $2.00$\,m. At the first violating step of each projected flow-based flight the
commanded position is already past the end of the plane, at $x=2.27$ to $2.52$\,m and $0.49$ to
$0.61$\,m ahead of the vehicle, and at every violating step the commanded position itself is clear of the
constraint. The plane is enforced over its span of $x$, and on this scene the span is checked at the
commanded position: once that position has left the span, the plan is released from the plane and
returns toward its route, and the vehicle, half a metre behind, follows it into the plane. The violating
steps under the tilt are thus the vehicle catching up with a plan that has already left the constraint,
which is why the budget makes them shallower without removing them. The baseline's projected flights leave
the corridor further to the side of their routes --- a median of $0.37$\,m, against $0.19$\,m for the
flow-based flights --- and stay there, and three to eight of ten of them fly the scene without a
violation.

Over the hump the unprojected flights stay level and pass through the roof. Under per-step projection at one
evaluation every flight climbs along the boundary and stops $0.12$ to $0.13$\,m short of the apex, at
$1.42$ to $1.43$\,m, where it stays until its $396$ control steps run out. At one and two evaluations
all three flow-based models stop this way, on $174$ of their $180$ flights, $0.01$ to $0.17$\,m short of
the apex; the other six clear the apex and then drift sideways across the boundary of the corridor wall.
Two properties of the setup meet at the apex. The plans never rise above $1.30$\,m, the top of the
altitude range of the demonstrations (dotted in \autoref{fig:uav-corridor-side}), while the apex asks for
$1.48$\,m. And each face of the roof, like the tilt, is enforced over its span of $x$ checked at the
commanded position, so while that position is short of the apex the rising face binds the whole plan,
including its part beyond the apex, where the face keeps rising. Both hold at every budget. At one and
two evaluations the corrected plan settles short of the apex; from three on it does not, and at twenty
the flights cross the apex grazing the boundary by about a centimetre and descend toward their route
altitude before the line. In altitude, projection does what the two constraints ask: the lowest altitude
under the tilt drops on $540$ of the $560$ projected flights, and the highest over the hump rises on all
$560$, while the unprojected flights fly level.

UAV-pillars is not drawn as a cost frontier: its result is the difference between a planner configuration's outcome on
D3IL-avoiding and the outcome of the same configuration flown, which is a comparison and not a trade-off
(\autoref{sec:res:uav:models}). UAV-s-curve is not drawn either: most of its flights end inverted
(\autoref{tab:uav-scurve}), and a frontier over its configurations would rank how early they fail, so that
scene gets the abort count printed beside each of its rows instead.

\FloatBarrier
\subsubsection{Projection Methods}\label{sec:res:uav:projection}
Per-step and endpoint projection separate on this scene only at twenty evaluations
(\autoref{tab:uav-corridor}; every selection rule in \autoref{app:uav-corridor-rules}). At three and five
evaluations endpoint projection is as cheap or cheaper, $44$ to $74$\,ms per action against $77$ to $94$ at
three evaluations and $73$ to $115$ against $112$ to $118$ at five, and it leaves more violating steps: at
three evaluations about two more per flight under the tilt ($6.9$ to $9.2$ against $5.0$ to $7.5$) and about
four more over the hump ($8.0$ to $10.4$ against $4.6$ to $6.5$), at five about one more. At twenty
evaluations endpoint projection leaves fewer and shallower violating steps than per-step projection under
both constraints: under the tilt $1.4$ to $2.5$ per flight, at most $1.3$\,cm deep, against $2.1$ to $3.2$
and $2.2$\,cm, and over the hump $0.2$ to $0.7$ against $1.8$ to $2.2$. The bar this thesis set for endpoint
projection, to overtake the per-step projector of \ac{DPCC} at a lower cost, is met on those two measures over
the hump at twenty evaluations with a single candidate, $0.2$ against $1.8$ violating steps per flight at $291$
against $355$\,ms per action, though on flights about thirty control steps longer, $316$ against $285$, so that
neither projector dominates the other; with four candidates it leaves $0.4$ to $0.7$ at $422$ to $425$\,ms. Under the tilt
its fewest violating steps come with four candidates at $417$\,ms, more time than the per-step projector's
$360$\,ms, and with a single candidate it is cheaper, $289$\,ms, but leaves $2.3$.

Endpoint projection has no row at $\nfe=1$ or $\nfe=2$ on this scene and will not get one: at the
activation threshold of $0.5$ the only active sampling step there is the terminal one, so there is no
guiding step for the method to act on (\autoref{sec:res:constraints:degenerate}). That is a property of
the method at those budgets, not a pending evaluation. The baseline has no endpoint rows for another
reason: endpoint projection adds its control to a velocity field (\autoref{sec:method:hardflow}), which
the diffusion model does not have.

The projection configuration of this scene differs from that of the other quadrotor scene
(\autoref{sec:setup:protocol:uav}), so the scenes are not pooled.

\subsubsection{Conclusion}\label{sec:res:uav:corridor:conclusion}

UAV-corridor shows two things the manipulation benchmarks could not. The first is that the concept of
\ac{DPCC}, a generative planner whose plans a projector holds inside the constraints, carries over to a new,
three-dimensional environment, against constraints that leave the horizontal plane and that no demonstration
satisfies. Projected, the flights descend under the tilt ($540$ of $560$) and climb over the hump (all $560$),
where the unprojected ones fly level through both, and the violating steps of a flight fall from $74.5$ to
$99.4$ (tilt) and $28.6$ to $36.6$ (hump) to $1.4$ to $10.4$ and $0.2$ to $8.3$ for the flow-based models and
to $0.8$ and $0.0$ for the baseline. What remains lies where each constraint ends, and the budget shrinks it.
Under the tilt it is the vehicle catching up with a plan the constraint has already released: the plane is
enforced over its span along the corridor at the commanded position, half a metre ahead of the vehicle, and
once that position has left the span the plan turns back toward its route while the vehicle is still inside.

The second is what the expert data decide. The demonstrations carry little information --- straight lines on
three lanes at one altitude per episode, a single path forward from any position (\autoref{tab:uav-demos})
--- and on them the generative model does not matter; the budget does. Without projection the flights of the
three flow-based objectives cannot be told apart; with it they leave the same violating steps at a shared
budget and projection method wherever the flights reach the end of the corridor, to within $1.6$ per flight
--- over the hump at one and two evaluations under per-step projection they do not reach it --- while under
the tilt the budget takes them from
about ten at one evaluation to one or two at twenty. The one place where projection falls short lies at the
edge of the data: per-step projection at one and two evaluations does not carry a plan over the apex of the
hump, at $1.48$\,m, above the $1.30$\,m that neither the plans nor the demonstrations exceed, and those flights
stop short of it. Against the baseline the flow-based models trade: a twenty-fourth (tilt) and a ninth (hump)
of its time per action at their cheapest, for more violating steps, at best $1.4$ and $0.2$ per flight against its
$0.8$ and $0.0$. The baseline's projected flights leave the corridor further to the side of their routes, a
median of $0.37$\,m against $0.19$\,m, and take more control steps to its end, $290$ and $321$ against $247$ to
$316$.

Between the two projection methods the scene shows a trade-off, and neither dominates. At twenty evaluations
endpoint projection leaves fewer violating steps under both constraints: over the hump with a single
candidate also at less time per action, $0.2$ against $1.8$ violating steps at $291$ against $355$\,ms, but on
flights about thirty control steps longer; under the tilt at more time, $1.4$ against $2.1$ at $417$ against
$360$\,ms. At three and five evaluations it is as cheap or cheaper but leaves as many violating steps or more,
and at one and two it has no guiding step.

\FloatBarrier
\subsection{UAV-s-curve}\label{sec:res:uav:scurve}

UAV-s-curve is read in the order of the two other scenes, with one step added. The generative models
fly the single route without projection first, each flow-based model at one, two and twenty network
evaluations and the diffusion baseline at the twenty it was trained with (\autoref{tab:uav-scurve}). The
configuration whose plans cross the finish line on the most flights is then the one the projector is
applied to (\autoref{tab:uav-scurve-projection}), and its plans, unprojected and projected, are what the
controller comparison of \autoref{sec:res:uav:controller} flies under the two tracking controllers.
Every cell is ten flights of the one route, one training seed, tracked by the cascaded geometric
controller of \autoref{sec:method:geometric} unless the row says otherwise. A flight ends in one of
three ways: it reaches the goal point of its route, it reaches the episode limit, or the divergence guard
ends it because the vehicle inverted (\autoref{sec:setup:metrics:uav}); crossing the finish line is latched
when it happens and does not end the flight. The last is the \emph{aborted} column,
and an aborted flight is never a success.

\subsubsection{Generative Models}\label{sec:res:uav:scurve:models}

\begin{table}[htpb]
  \caption[UAV-s-curve: the generative models without projection]{UAV-s-curve, every flow-based model at
  one, two and twenty network evaluations and the diffusion baseline at its own twenty, executed without
  projection under the cascaded geometric controller: ten flights of the single route per cell, one
  training seed. \emph{Success}: crosses the finish line in controlled flight (\autoref{sec:setup:metrics:uav});
  \emph{violating steps}: per flight, until
  the flight ends, and \emph{distance}, the final distance to the route's goal point, mean $\pm$ sample
  standard deviation over the flights; \emph{aborted}:
  ended by the divergence guard; \emph{ms/step}: time to compute one action. \selectedmark\ the
  configuration selected for projection; \baselinemark\ the baseline of record.}
  \label{tab:uav-scurve}
  \centering
  \small
  \begin{tabularx}{\linewidth}{@{}L r r r r r r@{}}
    \toprule
    Model & $\nfe$ & Success & \shortstack{Violating\\steps} & Distance [m] & Aborted & ms/step \\
    \midrule
    MeanFM        & 1  & 3/10 & $18.5 \pm 10.3$ & $2.023 \pm 1.193$ & 7/10  & 9.2 \\
    MeanFM        & 2  & 0/10 & $16.4 \pm 17.7$ & $2.515 \pm 0.796$ & 9/10  & 17.9 \\
    MeanFM        & 20 & 0/10 & $13.0 \pm 5.1$ & $2.772 \pm 0.264$ & 10/10 & 176.4 \\
    \addlinespace[3pt]
    CI-MeanFM, $\alpha_{\mathrm{end}}=0.2$ & 1  & 6/10 & $23.3 \pm 11.6$ & $1.306 \pm 1.319$ & 4/10 & 9.3 \\
    CI-MeanFM, $\alpha_{\mathrm{end}}=0.2$ & 2  & 1/10 & $21.3 \pm 22.3$ & $2.082 \pm 1.236$ & 7/10 & 18.0 \\
    CI-MeanFM, $\alpha_{\mathrm{end}}=0.2$ & 20 & 1/10 & $16.3 \pm 13.6$ & $2.511 \pm 0.843$ & 9/10 & 180.4 \\
    \addlinespace[3pt]
    \mainconfig{\selectedmark\ FM} & \mainconfig{1} & \mainconfig{9/10} & \mainconfig{$23.0 \pm 5.4$} & \mainconfig{$0.571 \pm 0.868$} & \mainconfig{1/10} & \mainconfig{8.9} \\
    FM            & 2  & 7/10 & $24.3 \pm 13.4$ & $1.110 \pm 1.328$ & 3/10 & 17.2 \\
    FM            & 20 & 6/10 & $26.1 \pm 20.0$ & $1.193 \pm 1.257$ & 4/10 & 168.3 \\
    \midrule
    \mainconfig{\baselinemark\ Diffusion} & \mainconfig{20} & \mainconfig{0/10} & \mainconfig{$119.4 \pm 199.4$} & \mainconfig{$2.432 \pm 0.674$} & \mainconfig{8/10} & \mainconfig{174.3} \\
    \bottomrule
  \end{tabularx}
\end{table}

No flight of any cell is free of violations: the flow-based flights violate on $13.0$ to $26.1$ control
steps on average, counted until the flight ends, and the baseline's on $119.4$, two of its flights ending on
the floor with close to five hundred. Instantaneous-velocity matching at one evaluation
crosses the finish line on nine flights of ten, one flight ending inverted, and ends $0.571$\,m from the
goal point on average; it is the configuration with the most crossings, followed by the same model at two
evaluations (seven) and at twenty and CI-MeanFM at one (six each). For every flow-based model the crossings
fall and the inversions rise from one to two to twenty evaluations: nine, seven and six crossings for
instantaneous-velocity matching, six, one and one for CI-MeanFM, three, none and none for analytic
average-velocity matching, whose ten flights at twenty evaluations all end inverted. The diffusion baseline
at its own twenty steps crosses on none, with eight inverted. The inversions come $11.6$ to $15.9$\,s into the
flight, at control steps $386$ to $530$, and at the second turn: $48$ of the $54$ flow-based inversions lie
between $x=0.2$ and $1.2$\,m, just past the crossover, where the route turns into its second straight. The
flights that cross reach the line at step $590$ to $635$. Every flight that crosses violates the constraints
on the way, and the reason is in the plans. On every flight of every configuration the commanded path cuts
the second inside corner of the crossover, passing $8$ to $19$\,cm inside its keep-out, where the
demonstrated route clears the constraints by $0.121$\,m, and that corner is where the violations lie. The plans
cross the gap about $0.2$\,m late: their commanded paths cross its middle at $x=0.21$ to $0.27$\,m, where the
Z-shaped route of the demonstrations crosses at $x=0$, half a metre from both corners. The
cascaded geometric controller follows the commanded path to within about a centimetre sideways on average,
and it never comes more than $1.4$\,cm closer to the corner than the command does; the $0.30$ to $0.33$\,m
between a flow-based flight and its commanded position is the lead of that position along the path, not a
sideways error.

The plans themselves show a second side of the budget (\autoref{fig:uav-scurve-plans}). At one and two
evaluations the four candidates of a plan scatter and zig-zag from one control step to the next, for both
models; at twenty they lie smooth along the route. The rougher plans are not the worse ones on this scene:
instantaneous-velocity matching crosses the finish line on nine flights with its roughest plans and on six
with its smoothest, and every flow-based model crosses most often at one evaluation. A plan of lower quality,
re-drawn at every control step, varies more from step to step, and on this route that variation appears to
help the controller through the turns rather than hinder it --- although, as the two panels of analytic
average-velocity matching show, it does not save every model's flights at the second corner.

\begin{figure}[htpb]
  \centering
  \includegraphics[width=\linewidth,height=0.70\textheight,keepaspectratio]{fig_uav_scurve_plans}
  \caption[UAV-s-curve: the plans without projection across the budget]{UAV-s-curve, the plans without
  projection, one flight per panel, drawn like \autoref{fig:raw-plans} but sampled: the plan the model
  emitted at every twentieth control step, all four candidates, each drawn as the measured planar position of
  its eight waypoints, green at the waypoint it is anchored at; the executed path is not drawn. Left, FM at
  $\nfe=1$, $2$ and $20$; right, MeanFM at $\nfe=1$ and $2$; the same flight of the route in every panel, the
  cascaded geometric controller, one training seed. A cross marks where the flight lost control.}
  \label{fig:uav-scurve-plans}
\end{figure}

\subsubsection{Projection}\label{sec:res:uav:scurve:projection}

The configuration \autoref{tab:uav-scurve} selects is instantaneous-velocity matching at one evaluation,
and it is flown under per-step projection on the tightened constraint set. Endpoint projection guides the
sample only on the sampling steps before the terminal one that the activation threshold selects
(\autoref{sec:setup:protocol:guiding}): at the default threshold of $0.5$ that is none at one evaluation, so
the method has nothing to act on at the selected budget and is not flown there.

\begin{table}[htpb]
  \caption[UAV-s-curve: the selected configuration under projection]{UAV-s-curve, FM at $\nfe=1$, the
  configuration selected in \autoref{tab:uav-scurve}, executed without projection and under per-step
  projection on the tightened constraint set, the cascaded geometric controller, ten flights per cell, one
  training seed. Columns as in \autoref{tab:uav-scurve}; \emph{ms/step} includes the projection. Per-step
  projection at random selection, the rule carried into \autoref{tab:uav-controller}: the three rules cross on
  five flights each, and random selection aborts least and violates least of the rules that abort least.
  Endpoint projection has no guiding step at $\nfe=1$.}
  \label{tab:uav-scurve-projection}
  \centering
  \small
  \begin{tabularx}{\linewidth}{@{}L L r r r r r@{}}
    \toprule
    Projection & Selection rule & Success & \shortstack{Violating\\steps} & Distance [m] & Aborted & ms/step \\
    \midrule
    none      & ---                   & 9/10 & $23.0 \pm 5.4$ & $0.571 \pm 0.868$ & 1/10 & 8.9 \\
    \addlinespace[3pt]
    per-step  & random                & 5/10 & $71.3 \pm 66.6$ & $0.572 \pm 0.345$ & 3/10 & 157.1 \\
    \addlinespace[3pt]
    endpoint  & ---                   & \multicolumn{5}{l}{--- no guiding step at $\nfe=1$} \\
    \bottomrule
  \end{tabularx}
\end{table}

Per-step projection does not keep the nine crossings. Under every selection rule the selected configuration
crosses on five flights of ten, the aborts rise from one to three or five, and the violating steps rise from
$23.0$ to between $37.2$ and $74.7$ per flight, while the time per control step rises from $8.9$ to between
$136$ and $164$\,ms, $126$ to $155$ of it the projection. Under random selection and cumulative projection cost
two flights each drop to the floor after crossing, which is where their spread comes from. The projector
does not reach what causes the violations: on this scene it binds the plan's measured positions, as \ac{DPCC}
does, and not the commanded position the vehicle is sent to, which UAV-corridor binds as well
(\autoref{sec:setup:protocol:uav}), and the projected commanded paths still cut the second inside corner,
$8$ to $9$\,cm inside its keep-out.
A per-step projector that also binds the commanded position, as on UAV-corridor, was not flown on this
scene.

\subsubsection{Caveat: The Tracking Controller}\label{sec:res:uav:controller}

On the manipulator, inverse kinematics and the joint-space controller lie between a generated plan and
the executed motion. The quadrotor is underactuated and open-loop unstable
(\autoref{sec:bg:quadrotor}), and a plan reaches the rotors only through the cascaded geometric
controller that tracks it (\autoref{sec:method:geometric}). The result of a flight is therefore a
property of the planner and the controller together, and UAV-s-curve, with its two sharp turns, is where
the choice of controller shows. Under the cascaded geometric controller most flights of analytic
average-velocity matching, of the consistency-interpolated model and of the baseline end inverted
(\autoref{tab:uav-scurve}), while instantaneous-velocity matching at one evaluation crosses the finish line
on nine flights of ten, each with $17$ to $37$ violating steps. What the controller changes is tested by
keeping the planner and its configuration and changing the controller, in closed loop: MuJoCo MPC (\autoref{sec:method:mjpc}) replaces the cascaded
geometric controller here and nowhere else.

The comparison is made on the configuration of \autoref{tab:uav-scurve-projection}, instantaneous-velocity
matching at one evaluation, with ten flights per cell (\autoref{tab:uav-controller},
\autoref{fig:uav-scurve-paths}): its plans, unprojected and under per-step projection with random selection,
under each controller. On the unprojected plans the two controllers succeed on the same nine flights of ten.
MuJoCo MPC removes the one inversion and brings every flight within $0.3$\,m of the goal point, a mean distance
of $0.297$ against $0.571$\,m, the tenth flight being scored unsafe for contact; but it follows the commanded
path less closely, $5$\,cm on average and up to $29$\,cm, the plans generated from its states cut the corner
deeper, and the violating steps rise from $23.0$ to $41.5$ per flight.

\begin{figure}[htpb]
  \centering
  \includegraphics[width=\linewidth]{fig_uav_scurve_paths}
  \caption[UAV-s-curve: the selected plans under the two tracking controllers]{UAV-s-curve from above, drawn
  like \autoref{fig:constraints-uav}: FM at $\nfe=1$, ten flights per panel; top, the unprojected plans;
  bottom, the plans under per-step projection (random selection, tightened constraint set); left, the
  cascaded geometric controller; right, MuJoCo MPC. Green: success; red: no success; a cross marks a flight
  that collided or that the divergence guard ended.}
  \label{fig:uav-scurve-paths}
\end{figure}

\begin{table}[htpb]
  \caption[UAV-s-curve: the same planner configuration under two tracking controllers]{UAV-s-curve, FM at
  $\nfe=1$, the configuration selected in \autoref{tab:uav-scurve}, flown in closed loop under each of the
  two tracking controllers,
  ten flights per cell, one training seed.
  \emph{Cascaded geometric} is the controller behind every other result in this chapter; \emph{MuJoCo
  MPC} (\autoref{sec:method:mjpc}) replaces it here and nowhere else. \emph{Per-step}: the per-step
  projector on the tightened constraint set with random selection, first of the rules that tie on crossings in
  \autoref{tab:uav-scurve-projection} by fewest aborts, then fewest violating steps. Columns as in
  \autoref{tab:uav-scurve}.}
  \label{tab:uav-controller}
  \centering
  \small
  \begin{tabularx}{\linewidth}{@{}L L r r r r@{}}
    \toprule
    Projection & Controller & Success & \shortstack{Violating\\steps} & Distance [m] & Aborted \\
    \midrule
    none      & cascaded geometric & 9/10 & $23.0 \pm 5.4$ & $0.571 \pm 0.868$ & 1/10 \\
    none      & MuJoCo MPC         & 9/10 & $41.5 \pm 18.4$ & $0.297 \pm 0.002$ & 0/10 \\
    \addlinespace[3pt]
    per-step  & cascaded geometric & 5/10 & $71.3 \pm 66.6$ & $0.572 \pm 0.345$ & 3/10 \\
    per-step  & MuJoCo MPC         & 0/10 & $105.5 \pm 46.0$ & $1.626 \pm 0.988$ & 7/10 \\
    \bottomrule
  \end{tabularx}
\end{table}
Ten flights per cell on the selected configuration. The gains are
held at their nominal values in every evaluation of this thesis and were not tuned; the two alternative
settings that exist scale the position loop by $\pm20\,\%$ and leave the attitude loop and the thrust limit
untouched. What the comparison supports is that the inversions are the controller's limit, not that no tuning
of this controller could fly the scene.

Under projection the cascaded geometric controller is the better of the two. It brings five of ten flights
across the finish line, and MuJoCo MPC none: seven
of its flights invert $15.7$ to $20.7$\,s into the flight, pressed against the outer wall of the second straight,
and the other three reach the goal region with more wall contact than the scene allows. MuJoCo MPC, which
keeps every unprojected flight upright, loses every projected one. No flight of the forty is free of violations under
either controller, because the corner cut is in the plans; MuJoCo MPC, following the projected plans less
closely, $12$\,cm on average against $8.5$\,cm for the cascaded law, violates on $105.5$ control steps per flight
against $71.3$.

The other reason MuJoCo MPC replaces the controller here and nowhere else is what it costs to run. The
time reported in every other table of this chapter is measured around the planner and stops there
(\autoref{sec:setup:metrics:avoiding}), so it prices no controller at all; pricing one means reading the
elapsed time of the flights themselves and dividing it by the control steps that executed
(\autoref{tab:uav-controller-cost}). On the selected plans without projection, where nothing lies between
the plan and the controller, a control step takes $14.0$\,ms under the cascaded geometric controller and
$135.8$\,ms under MuJoCo MPC. The planner's share of that is $8.9$ and $10.4$\,ms, so the difference belongs to
the controller: about $5$\,ms per step for the cascaded geometric law against about $125$\,ms for MuJoCo MPC,
which solves an optimisation problem at every control step where the other evaluates an algebraic law. MuJoCo
MPC is about twenty-four times the price of the cascaded geometric controller per control step, controller
and simulator step together on the unprojected pair, and by itself it costs twelve times as much per control
step as generating the plan --- some seven times what the operating
point of D3IL-avoiding spends on an entire action (\autoref{tab:avoiding-conclusion}).

\begin{table}[htpb]
  \caption[UAV-s-curve: what each tracking controller costs per control step]{UAV-s-curve, FM at
  $\nfe=1$ without projection, the configuration selected in \autoref{tab:uav-scurve}: the elapsed time of a
  flight divided by the control steps it executed, in milliseconds, under the two tracking controllers ---
  the one quantity in this chapter that contains the controller. \emph{The planner}: the planner time reported
  elsewhere; \emph{the controller and the simulator step}: the remainder. Ten flights per column, one training
  seed.}
  \label{tab:uav-controller-cost}
  \centering
  \small
  \begin{tabularx}{\linewidth}{@{}L r r@{}}
    \toprule
    Time per executed control step [ms] & Cascaded geometric & MuJoCo MPC \\
    \midrule
    The control loop                          & $14.0$ & $135.8$ \\
    \quad the planner                         & $8.9$ & $10.4$ \\
    \quad the controller and the simulator step & $5.2$ & $125.4$ \\
    \bottomrule
  \end{tabularx}
\end{table}
This is a decomposition of recorded elapsed times and not a measurement taken around the
controller call, which the evaluation does not instrument. The remainder row therefore carries the
simulator step with the controller, and its difference between the two columns is read as the controller's
cost, in two execution environments (MJX for MuJoCo MPC). One training seed, ten flights against ten.

What the comparison shows is that the controller matters. In this framework a plan reaches the rotors only
through the tracking controller, and the same projected planner configuration succeeds on five flights of ten or on none, and
the same unprojected configuration violates on $23$ control steps per flight or on $42$, depending on which
controller flies it in closed loop. On the configuration of record the cascaded geometric controller used throughout this thesis is
the better one: as many successes as MuJoCo MPC without projection, five against none with it, fewer violating
steps in both, at a twenty-fourth of the controller's cost per control step --- a tenth of the whole loop
unprojected, half of it projected. MuJoCo MPC's one gain, keeping the unprojected
flights upright and closer to the goal point, does not survive projection, and it is not adopted.

\FloatBarrier
\subsubsection{Conclusion}\label{sec:res:uav:scurve:conclusion}

UAV-s-curve shows that in this framework the controller matters. A plan reaches the rotors only through the
tracking controller, and on a route of two sharp turns the same planner configuration succeeds or fails by the controller that
flies it. On the configuration of record the cascaded geometric controller, the one every result of this
thesis is measured with, is the better of the two: without projection it succeeds on as many flights as
MuJoCo MPC, nine of ten, with fewer violating steps, and under projection on five flights where MuJoCo MPC
succeeds on none, at a twenty-fourth of MuJoCo MPC's controller cost per control step. MuJoCo MPC's one gain,
keeping
the unprojected flights upright and within $0.3$\,m of the goal point, does not survive projection. The
violations themselves are the plan's: every commanded path cuts the second inside corner of the crossover,
crossing the gap about $0.2$\,m later than the demonstrations do and passing $8$ to $19$\,cm inside the
corner's keep-out, and neither controller nor per-step projection, as configured, removes that cut.

Among the models instantaneous-velocity matching leads: at one evaluation it crosses the finish line on
nine flights of ten, the consistency-interpolated model on six and analytic average-velocity matching on
three, while most flights of the other two models and of the baseline invert at the second turn. That the
simplest objective leads is read here as a property of the data: the demonstrations are one route of straight segments
joined by arcs (\autoref{tab:uav-demos}), too simple to need what the average-velocity objectives are built
for. The scene adds a caveat on the budget: where a flight is easy to lose, more network evaluations are not
always better. For every flow-based model the crossings fall as the budget grows --- from nine to seven to
six for instantaneous-velocity matching, from six to one for the consistency-interpolated model and from
three to none for analytic average-velocity matching --- although the plans grow smoother with it
(\autoref{fig:uav-scurve-plans}): the rougher plans of one evaluation, re-drawn at every control step, vary
more from step to step, and on this route that variation appears to help the controller through the turns.
Instantaneous-velocity matching is sampled from initial noise of half its training scale
(\autoref{sec:setup:protocol:avoiding}); that bears on its order against the average-velocity models, not on
the controller comparison, which flies the same planner configuration under both controllers.

\section{Comparison across Environments}\label{sec:res:summary}

This section recaps the preceding ones and adds no new result. It answers, per environment, the first two
research questions (\autoref{sec:intro:rqs}): how the flow-based models compare with the diffusion baseline
and with one another, and whether endpoint projection improves on the per-step projection of \ac{DPCC}.

\subsection{Generative Models}\label{sec:res:summary:models}

The results follow the lines along which the study was designed. On the two D3IL tasks flow-based models
Pareto-dominate the diffusion model of \ac{DPCC} with its own temporal U-Net of $4.0$\,M parameters ---
CI-MeanFM and FM on D3IL-avoiding, MeanFM on D3IL-aligning; on the quadrotor they carry that result to a
new plant, and each scene adds what the plant or the data change.

\begin{description}
  \item[D3IL-avoiding: Pareto dominance over the baseline.] CI-MeanFM Pareto-dominates the baseline and MeanFM
    matches its saving at one episode in thirty. CI-MeanFM at one evaluation holds its success with constraint satisfaction of $1.000$ in $59.2$ control
    steps against $70.1$, at $18.1$ against $553.4$\,ms per action; MeanFM trades one episode in thirty
    ($0.967$) for the same saving; FM at one evaluation dominates the baseline too ($1.000$, $67.0$ steps,
    $17.3$\,ms). Against the baseline retrained at two denoising steps the saving is still twelvefold.
  \item[D3IL-aligning: the dominance under vision and contact.] On camera observations and contact, analytic average-velocity
    matching under per-step projection dominates the baseline's best projected row: nine violation-free
    contexts against four, a median box $0.179$ against $0.462$\,m from its target, $266$ against $809$\,ms.
    Without projection every non-dominated configuration is analytic average-velocity matching, which closes
    $84\,\%$ of the distance at twenty evaluations against $14$, $10$ and $-2\,\%$ for the others.
  \item[UAV-pillars: transfer of the D3IL-avoiding result to the quadrotor.] The Pareto-optimal operating points
    of D3IL-avoiding, flown again by the quadrotor: the declared constraints transfer exactly and the two
    average-velocity models keep their order, and the outcome departs from the table where the plant and its
    controller differ from the manipulator. The vehicle tracks more closely than the arm, and
    collisions with pillars outside the constraint sets, which end a quadrotor's flight, cost a third of the
    flights: success within the constraints $0.600$ to $0.700$.
  \item[UAV-corridor: projection in three dimensions and the role of the demonstrations.] The projection scheme
    of \ac{DPCC} carries over to three dimensions: the violating steps
    of a flight fall from $29$--$99$ to $0$--$10$. On its straight-line demonstrations the three objectives
    cannot be told apart and the budget decides; against the baseline the flow-based models trade a
    twenty-fourth (tilt) and a ninth (hump) of its time for more violating steps.
  \item[UAV-s-curve: the role of the tracking controller.] The same planner configuration succeeds or fails by the controller that
    flies it, and the cascaded geometric controller is the better one: as many successes as MuJoCo MPC without
    projection, five of ten against none with it, at a twenty-fourth of the controller's cost per control step.
    Instantaneous-velocity
    matching leads on this single, simple route, and more evaluations are not always better where a flight
    is easy to lose.
\end{description}

\subsection{Projection Methods}\label{sec:res:constraints}

Endpoint projection is the better projector on the constraint where the task needs a budget that gives it
steps to guide, as on D3IL-aligning; on UAV-corridor the two projectors trade against each other, and where one
or two evaluations solve the task endpoint projection has no step to guide.

\begin{description}
  \item[D3IL-avoiding (operated at $\nfe=1$).] Endpoint projection has no guiding step at the operating
    budget. Where it can act, with matched candidates, it is level or ahead on the constraint and cheaper on
    every model except instantaneous-velocity matching at three evaluations; on the average-velocity models
    at three it halves the time per control step.
  \item[D3IL-aligning ($\nfe=20$).] Endpoint projection is the better projector on the constraint: all ten
    contexts free of violations at the time per control step of per-step projection, and more contexts
    violation-free in seven of nine comparisons. Its advantage is on the constraint, not the task.
  \item[UAV-pillars ($\nfe=1$ and $2$).] Per-step projection only, at the D3IL-avoiding operating points: as
    there, it turns plans that violate the declared constraints into plans that violate none.
  \item[UAV-corridor ($\nfe=1$ to $20$).] A trade-off, with no clear dominance. At twenty evaluations endpoint
    projection leaves fewer violating steps under both constraints: over the hump with one candidate also at
    less time, $0.2$ against $1.8$ at $291$ against $355$\,ms, but on flights about thirty control steps
    longer; under the tilt at more time, $1.4$ against $2.1$ at $417$ against $360$\,ms. At three and five
    evaluations it leaves as many violating steps or more.
  \item[UAV-s-curve ($\nfe=1$).] Only per-step projection runs. It binds the measured position, leaves the
    corner cut of the plans in place and costs four of the nine crossings, at $126$ to $155$\,ms of
    projection per step.
\end{description}

\subsection{Conclusion}\label{sec:res:summary:conclusion}

The study is built in three steps, and its results follow them. The main claim is made and tested on
D3IL-avoiding, the benchmark of \ac{DPCC}, on which the endpoint-projection work was also evaluated: there
CI-MeanFM Pareto-dominates the diffusion model of \ac{DPCC} with its own temporal U-Net, MeanFM matches its
saving at one episode in thirty, and no
endpoint-projection configuration, which needs at least three evaluations to guide the sample, improves on their
operating point at one evaluation with per-step projection. D3IL-aligning extends the claim to a more complex
task, camera observations and an alignment through contact, and reaffirms it. The quadrotor scenes are mainly a
test of transfer: they carry the framework to a new plant and its tracking controller, and show where the result
holds and what the plant, the controller and the demonstrations change.

Within that structure the results form one picture. Where the planner learns from human
demonstrations, the average-velocity models lead: on D3IL-avoiding CI-MeanFM Pareto-dominates the diffusion model of
\ac{DPCC} with its own temporal U-Net, in fewer control steps at a thirtieth of its time per action, MeanFM
matches that saving at one episode in thirty, and on
D3IL-aligning analytic average-velocity matching does so again, closer to the target at a third of the time.
On these harder tasks the model and the projection matter together: per-step projection at one evaluation
completes D3IL-avoiding, and on D3IL-aligning, which needs twenty evaluations, endpoint projection helps as the
better projector on the constraint. The quadrotor scenes show the role of the demonstrations: generated for
this thesis and simple --- straight lanes on UAV-corridor, a single route on UAV-s-curve --- they leave the
objectives little to separate, so that instantaneous-velocity matching is enough; on UAV-corridor the three
objectives cannot be told apart, which is consistent with what the average-velocity objectives are built for,
and on
UAV-s-curve instantaneous-velocity matching leads. UAV-corridor also shows that the projection scheme of
\ac{DPCC} carries over to a three-dimensional environment against constraints no demonstration satisfies,
where the two projectors trade against each other without one dominating. And the quadrotor shows that the
plant and its controller matter: on UAV-pillars the vehicle, tracking more closely than the arm, flies into
the pillars the declared constraints leave out, and a collision ends its flight; on UAV-s-curve the same
plans succeed or fail by the controller that flies them, the cascaded geometric controller used throughout the
thesis being the better of the two.

Three findings cut across the environments.

\begin{enumerate}
  \item \textbf{The price of a projection is set by what the projector is handed.} A plan close to the
    feasible set is cheap to repair, which is why the flow-based models at one evaluation and endpoint
    projection's predicted endpoint are cheap, and the two-step diffusion baseline, handed a plan one step
    from noise, is not.
  \item \textbf{A task saturated at one evaluation hides what a task that rewards more evaluations exposes.}
    The self-referential target of the consistency-interpolated model, invisible on D3IL-avoiding, stops it
    improving on D3IL-aligning.
  \item \textbf{The demonstrations bound what a model can show.} On the straight lanes of UAV-corridor the
    objectives do not differ, and the plans stay within the altitudes the demonstrations cover, which is
    where per-step projection at the smallest budgets falls short.
\end{enumerate}

The combination each environment supports, and what it achieves:

\begin{description}
  \item[D3IL-avoiding.] MeanFM or CI-MeanFM at one evaluation, per-step projection, temporal consistency:
    success with constraint satisfaction of $0.967$ and $1.000$ in $58.6$ and $59.2$ control steps at $18.3$
    and $18.1$\,ms per action, against the baseline's $1.000$ in $70.1$ at $553.4$\,ms; five training seeds,
    the protocol of \ac{DPCC}.
  \item[D3IL-aligning.] MeanFM at twenty evaluations, endpoint projection, random selection: all ten contexts
    free of violations. The box ends in position, within the $1.8$\,cm of the task's own test, in none of
    them, but it comes close in two, $2.2$ and $4.7$\,cm from its target; in six more it closes $40$ to
    $75\,\%$ of the distance and ends $13$ to $27$\,cm away, and in the last two it does not move ($41$ and
    $45$\,cm). The median context ends $0.197$\,m from the target, a reduction of $57\,\%$ relative to the mean initial
    distance.
    The unprojected plan does not reach the tolerance either: its best two contexts end $2.8$ and $2.9$\,cm
    away.
  \item[UAV-pillars.] The D3IL-avoiding configurations, MeanFM or CI-MeanFM at one or two evaluations,
    per-step projection, temporal consistency: no violating step, and success within the constraints of
    $0.600$ to $0.700$, every failure a collision with a pillar outside the constraint sets.
  \item[UAV-corridor.] Any of the flow-based objectives tested at a budget --- the average-velocity models to
    three evaluations, instantaneous-velocity matching to twenty --- the budget chosen for the trade:
    violating steps per flight
    from $10.2$ at $27.4$\,ms (one evaluation, per-step) to $1.4$ at $416.9$\,ms (twenty, endpoint) under the
    tilt, and from $8.0$ at $71.9$\,ms (three, endpoint) to $0.2$ at $291.0$\,ms (twenty, endpoint, one
    candidate) over the hump; every flight reaches the end of the corridor.
  \item[UAV-s-curve.] FM at one evaluation without projection, under the cascaded geometric controller: nine
    of ten flights cross the finish line, none free of violations; MuJoCo MPC crosses as often without
    projection and on none with it, at twenty-four times the controller's cost per step.
\end{description}
